\chapter[Projective Toric Varieties and Polytopes]
{Projective Toric Varieties and\\ Polytopes}
\label{chap:projective-toric-varieties-polytopes}

\begin{chapteroverview}
Chapter~4 translated codimension-one geometry into data on the rays of
a fan. We now reverse the direction: the inequalities attached to a
divisor carve out a region in the character space \(M_\RR\), and the
lattice points in that region become global sections.

The central construction is
\[
D=\sum_\rho a_\rho D_\rho
\quad\longmapsto\quad
P_D
=
\left\{
 m\in M_\RR:
 \langle m,u_\rho\rangle\geq-a_\rho
 \text{ for every }\rho
\right\}.
\]
It begins a new chain of translations:
\[
\boxed{
\text{divisors}
\longrightarrow
\text{polyhedra and lattice points}
\longrightarrow
\text{linear systems and projective maps}.
}
\]

Later in the chapter, convexity and strict convexity of support
functions will detect basepoint freeness and ampleness. Normal fans
will then explain how polytopes build projective toric varieties.
\end{chapteroverview}

\begin{threedictionaries}
The three dictionaries now meet in the theory of sections.

\medskip

\begin{center}
\small
\begin{tabularx}{\linewidth}{@{}
  >{\raggedright\arraybackslash}X
  >{\raggedright\arraybackslash}X
  >{\raggedright\arraybackslash}X@{}}
\toprule
\textbf{Algebra}
&
\textbf{Geometry}
&
\textbf{Combinatorics}
\\
\midrule
characters and monomial sections
&
line bundles and linear systems
&
lattice points in \(P_D\)
\\
section multiplication
&
projective maps
&
Minkowski addition
\\
graded section rings
&
projective embeddings
&
polytopes and normal fans
\\
\bottomrule
\end{tabularx}
\end{center}
\end{threedictionaries}

\begin{surfacelaboratory}
For a toric surface, every inequality cuts out a half-plane in
\(M_\RR\). The polyhedron \(P_D\) can therefore be drawn explicitly.

Under a toric blowup, the new ray contributes one new inequality. On
the polytope side, this often cuts off a corner. We will repeatedly
use this visible operation to interpret vanishing conditions,
positivity, and projective maps on blown-up surfaces.
\end{surfacelaboratory}

% BEGIN INLINED: chapter05-section01-revised-v2.tex
\section{Divisor polyhedra and global sections}
\label{sec:divisor-polyhedra-global-sections}
\index{divisor polyhedron}
\index{global section!monomial basis}
\index{lattice point!as a section}
\index{corner cut!blowup}

\begin{agcheatsheet}
Let \(X\) be a normal irreducible variety and let \(D\) be a Weil
divisor. The sheaf \(\mathcal O_X(D)\) is defined inside the constant
sheaf of rational functions by
\[
\Gamma(U,\mathcal O_X(D))
=
\left\{
 f\in\Bbbk(X)^\times:
 \bigl(\operatorname{div}(f)+D\bigr)|_U\geq0
\right\}
\cup\{0\}.
\]
Thus
\[
H^0(X,\mathcal O_X(D))
=
\left\{
 f\in\Bbbk(X)^\times:
 \operatorname{div}(f)+D\geq0
\right\}
\cup\{0\}.
\]

The inequality has a concrete meaning. The rational function \(f\)
may have poles, but along every prime divisor its pole order must be
no larger than the multiplicity supplied by \(D\). Equivalently,
after adding the zeros and poles of \(f\) to \(D\), the result must be
effective.

When \(D\) is Cartier, \(\mathcal O_X(D)\) is a line bundle and has a
canonical rational section \(s_D\) with
\[
\operatorname{div}(s_D)=D.
\]
Multiplication by this section identifies rational functions with
rational sections:
\[
f\longmapsto f s_D.
\]
Under this identification, \(f s_D\) is regular exactly when
\[
\operatorname{div}(f)+D\geq0.
\]
Hence the inequality is precisely the regularity condition for the
corresponding section of the line bundle.

For a non-Cartier Weil divisor on a normal variety, the same formula
defines the associated rank-one reflexive sheaf.
\end{agcheatsheet}

Now let
\[
D=\sum_{\rho\in\Sigma(1)}a_\rho D_\rho
\]
be torus invariant. By the principal-divisor formula from
\cref{thm:principal-divisor-character},
\[
\operatorname{div}(\character{m})+D
=
\sum_{\rho\in\Sigma(1)}
\bigl(\langle m,u_\rho\rangle+a_\rho\bigr)D_\rho.
\]
Therefore
\[
\character{m}\in H^0(X_\Sigma,\mathcal O_{X_\Sigma}(D))
\]
precisely when
\[
\langle m,u_\rho\rangle\geq-a_\rho
\qquad
\text{for every }\rho\in\Sigma(1).
\]
When \(D\) is Cartier, this says equivalently that the rational
section \(\character{m}s_D\) is regular. These inequalities cut out
a polyhedron in \(M_\RR\).


\subsection{The monomial basis theorem}

\begin{proposition}[Global sections decompose into characters]
\label{prop:global-sections-characters}
Let
\[
D=\sum_{\rho\in\Sigma(1)}a_\rho D_\rho
\]
be a torus-invariant Weil divisor on \(X_\Sigma\). Then
\[
H^0(X_\Sigma,\mathcal O_{X_\Sigma}(D))
=
\bigoplus_{\operatorname{div}(\character{m})+D\geq0}
\Bbbk\,\character{m}.
\]
\end{proposition}

\noindent
\textit{Reference.}
This is Cox--Little--Schenck, Proposition~4.3.2.

\begin{proof}
A global section \(f\) restricts to a regular function on the dense
torus, because \(D\) has no support there. Hence
\[
f|_T\in\Bbbk[M].
\]
The space of global sections is stable under the torus action, while
\(\Bbbk[M]\) has the character basis
\[
\{\character{m}:m\in M\}.
\]
Therefore the space of sections is the direct sum of the character
weight spaces that it contains. Finally,
\[
\character{m}\in H^0(X_\Sigma,\mathcal O_{X_\Sigma}(D))
\]
if and only if
\[
\operatorname{div}(\character{m})+D\geq0.
\]
\end{proof}

\begin{observation}
The short representation-theoretic phrase ``weight-space
decomposition'' is optional language for a simple fact: Laurent
monomials form a basis of \(\Bbbk[M]\), and a torus-stable subspace is
spanned by the monomials it contains.
\end{observation}

\begin{definition}[The polyhedron of a divisor]
\label{def:divisor-polyhedron}
For a torus-invariant Weil divisor
\[
D=\sum_{\rho\in\Sigma(1)}a_\rho D_\rho,
\]
define
\[
\boxed{
P_D
=
\left\{
 m\in M_\RR:
 \langle m,u_\rho\rangle\geq-a_\rho
 \text{ for every }\rho\in\Sigma(1)
\right\}.
}
\]
\end{definition}

Each ray contributes one closed half-space. The coefficient \(a_\rho\)
controls the position of its boundary hyperplane:
\[
\langle m,u_\rho\rangle=-a_\rho.
\]
Increasing \(a_\rho\) moves that wall outward and weakens the pole
restriction along \(D_\rho\).

\begin{theorem}[Lattice points are monomial sections]
\label{thm:lattice-points-global-sections}
For every torus-invariant Weil divisor \(D\),
\[
\boxed{
H^0(X_\Sigma,\mathcal O_{X_\Sigma}(D))
=
\bigoplus_{m\in P_D\cap M}
\Bbbk\,\character{m}.
}
\]
In particular, whenever \(P_D\cap M\) is finite,
\[
h^0(X_\Sigma,\mathcal O_{X_\Sigma}(D))
=
|P_D\cap M|.
\]
\end{theorem}

\noindent
\textit{Reference.}
This is Cox--Little--Schenck, Proposition~4.3.3.

\begin{proof}
By \cref{thm:principal-divisor-character},
\[
\operatorname{div}(\character{m})+D\geq0
\]
if and only if
\[
\langle m,u_\rho\rangle+a_\rho\geq0
\]
for every ray \(\rho\). This is exactly the condition
\[
m\in P_D\cap M.
\]
Now apply \cref{prop:global-sections-characters}.
\end{proof}

The theorem is the first major bridge of this chapter:
\[
\boxed{
\text{global sections}
\quad\longleftrightarrow\quad
\text{lattice points in a polyhedron}.
}
\]

\begin{example}[The simplex of sections on \(\PP^2\)]

Let
\[
u_0=-e_1-e_2,
\qquad
u_1=e_1,
\qquad
u_2=e_2,
\]
and consider
\[
D=dD_0,
\qquad
d\geq0.
\]
For \(m=(a,b)\in M\simeq\ZZ^2\), the inequalities defining \(P_D\)
are
\[
a\geq0,
\qquad
b\geq0,
\qquad
a+b\leq d.
\]
Hence
\[
P_D=d\Delta_2
=
\Conv\bigl((0,0),(d,0),(0,d)\bigr).
\]

\begin{figure}[ht]
\centering
\begin{tikzpicture}[scale=0.88]
  % Fan
  \begin{scope}[xshift=-4.6cm]
    \fill[DictionaryBlue!9]
      (0,0)--(2.5,0)--(0,2.5)--cycle;
    \fill[SurfaceGreen!9]
      (0,0)--(0,2.5)--(-1.8,-1.8)--cycle;
    \fill[ExampleGold!10]
      (0,0)--(-1.8,-1.8)--(2.5,0)--cycle;

    \draw[very thick,DictionaryBlue,-{Latex[length=2.3mm]}]
      (0,0)--(2.85,0);
    \draw[very thick,SurfaceGreen,-{Latex[length=2.3mm]}]
      (0,0)--(0,2.85);
    \draw[very thick,ExampleGold,-{Latex[length=2.3mm]}]
      (0,0)--(-2.1,-2.1);

    \node[DictionaryBlue,below] at (2.6,0) {\(u_1=e_1\)};
    \node[SurfaceGreen,left] at (0,2.6) {\(u_2=e_2\)};
    \node[ExampleGold,above left] at (-1.9,-1.9)
      {\(u_0=-e_1-e_2\)};

    \node[fill=DictionaryBlue!12,rounded corners,inner sep=2pt]
      at (1.0,0.82) {\(a\geq0\)};
    \node[fill=SurfaceGreen!12,rounded corners,inner sep=2pt]
      at (-0.55,0.85) {\(b\geq0\)};
    \node[fill=ExampleGold!14,rounded corners,inner sep=2pt]
      at (0.78,-0.55) {\(a+b\leq d\)};

    \node[font=\bfseries] at (0.25,3.55)
      {rays become inequalities};
  \end{scope}

  % Polytope d=3
  \begin{scope}[xshift=4.55cm,scale=0.95]
    \draw[->,black!65] (-0.25,0)--(3.7,0)
      node[below] {\(a\)};
    \draw[->,black!65] (0,-0.25)--(0,3.7)
      node[left] {\(b\)};

    \filldraw[
      fill=PredictionTeal!13,
      draw=PredictionTeal,
      very thick
    ] (0,0)--(3,0)--(0,3)--cycle;

    \foreach \a in {0,...,3} {
      \foreach \b in {0,...,3} {
        \pgfmathtruncatemacro{\s}{\a+\b}
        \ifnum\s<4
          \fill (\a,\b) circle (2.2pt);
        \fi
      }
    }

    \node[below left] at (0,0) {\((0,0)\)};
    \node[below] at (3,0) {\((d,0)\)};
    \node[left] at (0,3) {\((0,d)\)};
    \node[rotate=-45,PredictionTeal] at (1.85,1.55)
      {\(a+b=d\)};
    \node[font=\bfseries] at (1.45,4.05)
      {\(P_{dD_0}=d\Delta_2\)};
    \node[align=center] at (1.55,-0.85)
      {lattice points for \(d=3\)};
  \end{scope}
\end{tikzpicture}
\caption{For \(D=dD_0\) on \(\PP^2\), each ray contributes one
inequality. Their intersection is the dilated standard simplex, and
its lattice points index the degree-\(d\) monomial sections.}
\label{fig:p2-divisor-polytope-sections}
\end{figure}

The lattice point \((a,b)\) corresponds to
\[
\character{(a,b)}
=
\left(\frac{z_1}{z_0}\right)^a
\left(\frac{z_2}{z_0}\right)^b.
\]
As a section of \(\mathcal O_{\PP^2}(dD_0)\), this becomes the
homogeneous monomial
\[
z_0^{d-a-b}z_1^az_2^b.
\]
Thus
\[
H^0(\PP^2,\mathcal O_{\PP^2}(d))
\]
has the familiar basis of homogeneous degree-\(d\) monomials.

\end{example}

\begin{pauseandthink}
Count the lattice points in \(d\Delta_2\) by grouping them according
to the value of \(a+b\). Recover
\[
h^0(\PP^2,\mathcal O_{\PP^2}(d))
=
\binom{d+2}{2}.
\]
\end{pauseandthink}

\subsection{Polyhedra need not be bounded}

\begin{example}[The structure sheaf of \(\AAA^2\)]
For the first-quadrant cone
\[
\sigma=\Cone{e_1,e_2},
\]
take \(D=0\). Then
\[
P_0
=
\left\{(a,b)\in\RR^2:a\geq0,\ b\geq0\right\}
=
\sigma^\vee.
\]
This polyhedron is unbounded, and
\[
P_0\cap M=\ZZ_{\geq0}^2.
\]
The lattice-point theorem gives
\[
H^0(\AAA^2,\mathcal O_{\AAA^2})
=
\bigoplus_{a,b\geq0}\Bbbk x^ay^b
=
\Bbbk[x,y].
\]
The infinitely many lattice points reflect the infinitely many
regular functions on the affine plane.
\end{example}

For a general divisor, the recession cone of \(P_D\) is independent
of the coefficients of \(D\):
\[
\operatorname{rec}(P_D)
=
\left\{
 v\in M_\RR:
 \langle v,u_\rho\rangle\geq0
 \text{ for every }\rho
\right\}.
\]
It records the directions in which the polyhedron can continue
forever.

\begin{proposition}[Complete fans give bounded section polyhedra]
\label{prop:complete-fan-bounded-PD}
If \(\Sigma\) is complete, then for every torus-invariant Weil divisor
\(D\):
\begin{enumerate}[label=(\alph*),wide=0pt]
\item \(P_D\) is a polytope, possibly empty;
\item \(H^0(X_\Sigma,\mathcal O_{X_\Sigma}(D))\) is finite
dimensional;
\item \(H^0(X_\Sigma,\mathcal O_{X_\Sigma})=\Bbbk\).
\end{enumerate}
\end{proposition}

\noindent
\textit{Reference.}
This is Cox--Little--Schenck, Proposition~4.3.8.

\begin{proof}[Conceptual proof sketch]
A direction \(v\) lies in the recession cone of \(P_D\) precisely when
\[
\langle v,u_\rho\rangle\geq0
\]
for every ray. If the fan is complete, its cones fill all of
\(N_\RR\), so the only linear functional nonnegative in every
direction is \(v=0\). Hence \(P_D\) has no nonzero unbounded
direction and is therefore bounded.

A bounded polyhedron has finitely many lattice points, giving part
\textup{(b)}. For \(D=0\), completeness forces
\[
P_0=\{0\},
\]
so the only global regular functions are constants.
\end{proof}

\subsection{Blowing up cuts off a corner}

Let
\[
\pi:X=\Blp{\PP^2}{p}\longrightarrow\PP^2,
\qquad
p=[1:0:0],
\]
be the toric blowup. The fan has rays
\[
u_0=-e_1-e_2,
\qquad
u_1=e_1,
\qquad
u_E=e_1+e_2,
\qquad
u_2=e_2.
\]
Recall that
\[
[D_0]=H,
\qquad
[E]=E.
\]
For integers \(0\leq r\leq d\), choose the invariant representative
\[
D=dD_0-rE
\]
of the class \(dH-rE\). Its polytope is cut out by
\[
a\geq0,
\qquad
b\geq0,
\qquad
r\leq a+b\leq d.
\]
The new ray \(u_E=e_1+e_2\) contributes the new inequality
\[
a+b\geq r.
\]

\begin{prediction}
Start with the simplex \(d\Delta_2\), whose lattice points represent
all degree-\(d\) monomials in \(z_0,z_1,z_2\).

Before looking at the next figure, predict which corner should be
removed when we require the corresponding polynomial to vanish at
\[
p=[1:0:0]
\]
to order at least \(r\). Which exponent controls that order of
vanishing?
\end{prediction}

\begin{figure}[ht]
\centering
\begin{tikzpicture}[scale=0.92]
  % Blowup fan
  \begin{scope}[xshift=-4.6cm]
    \fill[DictionaryBlue!8]
      (0,0)--(2.5,0)--(1.25,1.25)--cycle;
    \fill[PredictionTeal!10]
      (0,0)--(1.25,1.25)--(0,2.5)--cycle;
    \fill[SurfaceGreen!8]
      (0,0)--(0,2.5)--(-1.8,-1.8)--cycle;
    \fill[ExampleGold!9]
      (0,0)--(-1.8,-1.8)--(2.5,0)--cycle;

    \draw[very thick,DictionaryBlue,-{Latex[length=2.3mm]}]
      (0,0)--(2.85,0);
    \draw[very thick,PredictionTeal,-{Latex[length=2.3mm]}]
      (0,0)--(2.1,2.1);
    \draw[very thick,SurfaceGreen,-{Latex[length=2.3mm]}]
      (0,0)--(0,2.85);
    \draw[very thick,ExampleGold,-{Latex[length=2.3mm]}]
      (0,0)--(-2.1,-2.1);

    \node[DictionaryBlue,below] at (2.62,0) {\(u_1\)};
    \node[PredictionTeal,above right] at (1.88,1.88)
      {\(u_E=e_1+e_2\)};
    \node[SurfaceGreen,left] at (0,2.62) {\(u_2\)};
    \node[ExampleGold,above left] at (-1.9,-1.9) {\(u_0\)};

    \node[fill=PredictionTeal!14,rounded corners,inner sep=2pt]
      at (1.0,0.72) {new wall: \(a+b\geq r\)};
    \node[font=\bfseries] at (0.3,3.55)
      {star subdivision};
  \end{scope}

  % Truncated triangle d=3 r=1
  \begin{scope}[xshift=4.55cm,scale=0.95]
    \draw[->,black!65] (-0.25,0)--(3.7,0)
      node[below] {\(a\)};
    \draw[->,black!65] (0,-0.25)--(0,3.7)
      node[left] {\(b\)};

    \fill[WarningRed!10]
      (0,0)--(1,0)--(0,1)--cycle;
    \draw[WarningRed,dashed,very thick]
      (1,0)--(0,1);

    \filldraw[
      fill=PredictionTeal!13,
      draw=PredictionTeal,
      very thick
    ] (1,0)--(3,0)--(0,3)--(0,1)--cycle;

    \foreach \a in {0,...,3} {
      \foreach \b in {0,...,3} {
        \pgfmathtruncatemacro{\s}{\a+\b}
        \ifnum\s<4
          \ifnum\s>0
            \fill (\a,\b) circle (2.2pt);
          \fi
        \fi
      }
    }

    \draw[WarningRed,very thick] (0,0) circle (3pt);
    \node[WarningRed,below left] at (0,0)
      {removed};
    \node[PredictionTeal,rotate=-45] at (0.58,0.72)
      {\(a+b=r\)};
    \node[PredictionTeal,rotate=-45] at (2.0,1.45)
      {\(a+b=d\)};
    \node[font=\bfseries] at (1.5,4.05)
      {\(P_{dD_0-rE}\)};
    \node[align=center] at (1.55,-0.85)
      {shown for \(d=3\), \(r=1\)};
  \end{scope}
\end{tikzpicture}
\caption{Blowing up the fixed point inserts the ray \(e_1+e_2\),
and the coefficient \(-r\) of the exceptional divisor inserts the
inequality \(a+b\geq r\). In the section polytope, this cuts off the
corner corresponding to monomials that do not vanish sufficiently at
the blown-up point.}
\label{fig:blowup-truncated-section-polytope}
\end{figure}

The lattice point \((a,b)\) corresponds to
\[
z_0^{d-a-b}z_1^az_2^b.
\]
At the point \(p=[1:0:0]\), its order of vanishing is
\[
a+b.
\]
Consequently,
\[
H^0\bigl(X,\mathcal O_X(dH-rE)\bigr)
\]
has the monomial basis indexed by
\[
\left\{
(a,b)\in\ZZ_{\geq0}^2:
 r\leq a+b\leq d
\right\}.
\]
The exceptional inequality is therefore the algebraic condition of
vanishing to order at least \(r\) at the blown-up point.

\begin{surfacelaboratory}
For \(0\leq r\leq d\), use the truncated simplex to answer the
following.

\begin{enumerate}[label=(\alph*),wide=0pt]
\item Compute
\[
h^0\bigl(X,\mathcal O_X(dH-rE)\bigr)
\]
by counting lattice points.

\item Check that your answer equals
\[
\binom{d+2}{2}-\binom{r+1}{2}.
\]

\item Interpret the subtracted term as the number of independent
conditions imposed by vanishing to order at least \(r\) at one point.

\item Draw the polytopes for \(d=4\) and \(r=0,1,2,3,4\). What
happens when \(r=d\)?
\end{enumerate}
\end{surfacelaboratory}

\begin{guidingquestion}
The polyhedron \(P_D\) tells us all global sections of
\(\mathcal O_X(D)\). How can its shape tell whether, at every point of
\(X_\Sigma\), at least one of those sections is nonzero?

Equivalently, what condition on the support function \(\psi_D\) makes
\(\mathcal O_X(D)\) generated by its global sections?
\end{guidingquestion}
% END INLINED: chapter05-section01-revised-v2.tex

% BEGIN INLINED: chapter05-section02-projective-maps-v2.tex
\section{Projective maps from lattice points}
\label{sec:projective-maps-lattice-points}
\index{projective toric variety}
\index{projective map!from lattice points}
\index{very ample divisor}
\index{ample divisor}

\begin{agcheatsheet}
Let \(\mathcal L\) be a line bundle on a complete variety \(X\), and
let
\[
V\subseteq H^0(X,\mathcal L)
\]
be a finite-dimensional space of global sections.

If the sections in \(V\) have no common zero, then they define a
morphism
\[
\phi_V:X\longrightarrow \PP(V^\vee).
\]
After choosing a basis \(s_0,\dots,s_r\) of \(V\), this map is given
in a local trivialization of \(\mathcal L\) by
\[
x\longmapsto [s_0(x):\cdots:s_r(x)].
\]
Changing the trivialization multiplies all coordinates by the same
nonzero scalar, so the point of projective space is well defined.

When
\[
V=H^0(X,\mathcal L),
\]
this is the map associated to the \emph{complete linear system}.

The line bundle \(\mathcal L\) is \emph{very ample} if its complete
linear system gives a closed embedding into projective space. It is
\emph{ample} if some positive tensor power
\[
\mathcal L^{\otimes k}
\]
is very ample.
\end{agcheatsheet}

% BEGIN CLASSROOM ADDITION: notes/20261014.tex (evaluation and ratios).
\begin{observation}[Evaluation in a one-dimensional fiber]
\label{obs:classroom-evaluation-ratios}
Intrinsically, evaluation at \(x\) is a linear map
\[
\operatorname{ev}_x:V\longrightarrow\mathcal L|_x.
\]
Its target is the fiber, rather than a canonically chosen copy of
\(\Bbbk\). When evaluation is nonzero, identifying that fiber with
\(\Bbbk\) gives a nonzero element of \(V^\vee\), well defined up to
scalar. This explains the point of \(\PP(V^\vee)\) without choosing
coordinates on the fiber.

On the open set where \(s_j\ne0\), the affine coordinates of the
projective map are the regular ratios \(s_i/s_j\). Thus division
requires a nonvanishing denominator. The basepoints of a chosen
subspace \(V\) concern that subspace; a line bundle is basepoint free
when its entire space of global sections has no common zero.
\end{observation}
% END CLASSROOM ADDITION

Let \(X_\Sigma\) be a complete toric variety, and let
\[
D=\sum_{\rho\in\Sigma(1)} a_\rho D_\rho
\]
be a torus-invariant Cartier divisor. By \cref{thm:lattice-points-global-sections},
its monomial sections are indexed by the lattice points of
\[
P_D=
\left\{
 m\in M_\RR:
 \langle m,u_\rho\rangle\geq -a_\rho
 \text{ for every }\rho
\right\}.
\]
Write
\[
P_D\cap M=\{m_0,\dots,m_r\}.
\]

Each lattice point determines a global section
\[
s_i=\character{m_i}s_D,
\]
where \(s_D\) is the distinguished rational section from
\cref{sec:divisor-polyhedra-global-sections}.

\subsection{From sections to a projective map}

If the sections \(s_0,\dots,s_r\) have no common zero, then they give
a morphism
\[
\phi_D:X_\Sigma\longrightarrow \PP^r.
\]
On the dense torus, the section \(s_D\) is nowhere zero, so it can be
cancelled from all homogeneous coordinates. Therefore
\[
\phi_D|_T(t)=
[\character{m_0}(t):\cdots:\character{m_r}(t)].
\]
This is the most concrete bridge from lattice points to projective
geometry: \emph{the exponents become homogeneous coordinates on the
image}. 

\begin{observation}
If the whole polytope is translated by one lattice vector
\(m\in M\), then every exponent becomes \(m_i+m\). On the torus, this
multiplies every projective coordinate by the common factor
\(\character{m}(t)\), so the projective map does not change.

Thus the projective map depends only on the affine shape of the set of
lattice points, not on its absolute position in \(M\).
\end{observation}

\begin{definition}
When \(D\) is basepoint free, we write
\[
\phi_D:X_\Sigma\longrightarrow \PP^{|P_D\cap M|-1}
\]
for the morphism defined by the complete linear system
\(|D|\).
\end{definition}

At this stage, the word ``basepoint free'' is just a condition saying
that \(\phi_D\) is everywhere defined. The next section will explain
how to detect that condition directly from the support function.

\subsection{Any finite set of lattice points gives a projective toric variety}

The same construction makes sense before a divisor enters the story.
Let
\[
A=\{m_0,\dots,m_r\}\subseteq M
\]
be any finite set of lattice points. Define a map on the torus by
\[
\phi_A:T\longrightarrow \PP^r,
\qquad
\phi_A(t)=
[\character{m_0}(t):\cdots:\character{m_r}(t)].
\]

\begin{proposition}[Projective toric varieties from lattice points]
\label{prop:projective-toric-variety-from-lattice-points}
Let
\[
Y_A=\overline{\phi_A(T)}\subseteq \PP^r.
\]
Then \(Y_A\) is a projective toric variety.
\end{proposition}

\noindent
\textit{Reference.}
This is the construction in Cox--Little--Schenck, \S7.1.

\begin{proof}[Conceptual proof sketch]
Let the torus act on \(\PP^r\) diagonally by
\[
t\cdot[z_0:\cdots:z_r]
=
[\character{m_0}(t)z_0:\cdots:\character{m_r}(t)z_r].
\]
Then \(\phi_A(T)\) is precisely the orbit of the point
\([
1:\cdots:1
]\).
Its closure is therefore a projective variety containing a dense torus
orbit, so it is a projective toric variety.
\end{proof}

This proposition says that \emph{projective toric varieties are born
from finite sets of lattice points}. When \(A=P_D\cap M\), the image
of \(\phi_D\) is exactly one such variety.

\subsection{The first examples}

\begin{example}[\(\PP^2\) and the Veronese picture]
Let \(H=D_0\) on \(\PP^2\). From the \(\PP^2\) example in the previous section, the polytope of \(dH\) is
\[
P_{dH}=d\Delta_2,
\]
whose lattice points are the pairs \((a,b)\) with
\[
a,b\geq 0,
\qquad
a+b\leq d.
\]
The corresponding sections are the degree-\(d\) monomials
\[
x_1^ax_2^bx_0^{d-a-b},
\]
more conveniently written as
\[
x_0^{d-a-b}x_1^ax_2^b.
\]
For \(d=1\), the three lattice points of \(\Delta_2\) give the three
homogeneous coordinates, and
\[
\phi_H:\PP^2\longrightarrow \PP^2
\]
is the identity map.

For general \(d\), the lattice points of \(d\Delta_2\) give the
\(d\)-th Veronese map
\[
\PP^2\longrightarrow \PP^{\binom{d+2}{2}-1}.
\]
Thus dilating the triangle produces the familiar projective embedding
by degree-\(d\) plane curves.
\end{example}

% BEGIN CLASSROOM ADDITION: notes/20261014.tex (the square and Segre map).
\begin{example}[The unit square and the Segre quadric]
\label{ex:classroom-segre-square}
On \(\PP^1\times\PP^1\), take the rays
\(e_1,e_2,-e_1,-e_2\), with divisor coefficients \(0,0,1,1\),
respectively. The associated line bundle is \(\mathcal O(1,1)\), and
\[
P_D=[0,1]^2.
\]
The four lattice points give the torus map
\[
(s,t)\longmapsto[1:s:t:st].
\]
In homogeneous coordinates this extends to
\[
\begin{aligned}
\PP^1\times\PP^1&\longrightarrow\PP^3,\\
([u_0:u_1],[v_0:v_1])
&\longmapsto[u_0v_0:u_1v_0:u_0v_1:u_1v_1].
\end{aligned}
\]
The image is the quadric
\[
z_0z_3=z_1z_2.
\]
Indeed, its coordinates form a nonzero rank-one matrix
\(\left(\begin{smallmatrix}z_0&z_2\\z_1&z_3\end{smallmatrix}\right)\).
Such a matrix recovers its two projective factors uniquely, and this
recovery is regular on the standard coordinate charts. Thus the map
is the Segre embedding. The square therefore gives both a familiar
surface and a concrete equation for its projective image.
\end{example}
% END CLASSROOM ADDITION

\begin{example}[A blowup map from a truncated triangle]
On
\[
X=\Blp{\PP^2}{p},
\]
consider
\[
D=2H-E.
\]
From \cref{fig:blowup-truncated-section-polytope},
\[
P_D=
\left\{
(a,b)\in\RR^2:
 a,b\geq 0,
 1\leq a+b\leq 2
\right\}.
\]
Its lattice points are
\[
(1,0),\ (2,0),\ (1,1),\ (0,2),\ (0,1).
\]
Hence the complete linear system of \(2H-E\) gives a map
\[
\phi_{2H-E}:X\longrightarrow \PP^4.
\]
On the torus, it is the monomial map
\[
(t_1,t_2)
\longmapsto
[t_1:t_1^2:t_1t_2:t_2^2:t_2].
\]
Compared with the full triangle for \(2H\), the missing lattice point
\((0,0)\) reflects the vanishing condition imposed at the blown-up
point.
\end{example}

% BEGIN CLASSROOM ADDITION: notes/20261014.tex (projection and its resolution).
% The lecture uses [0:0:1]; coordinates are permuted to the chapter's [1:0:0].
\begin{example}[Resolving projection from a point]
\label{ex:classroom-projection-pencil}
Keep the chapter's coordinates \(p=[1:0:0]\), and consider the
two-dimensional subsystem
\[
V=\langle z_1,z_2\rangle\subset H^0(\PP^2,\mathcal O(1)).
\]
It defines the rational projection
\[
\phi_V:\PP^2\dashrightarrow\PP^1,\qquad
[z_0:z_1:z_2]\longmapsto[z_1:z_2],
\]
whose basepoint is \(p\). Geometrically, it records the line through
\(p\) and the chosen point. The segment with endpoints \(e_1,e_2\)
indexes this subsystem of \(|H|\), rather than its complete set of
sections on \(\PP^2\).

On the blowup \(X=\Blp{\PP^2}{p}\), removing the common exceptional
factor from the pullbacks of \(z_1,z_2\) gives two generating sections
of \(H-E\). Here the complete section polytope is exactly
\[
P_{H-E}=\{(a,b):a,b\geq0,\ a+b=1\}.
\]
The rational projection has become a morphism:
\[
\begin{tikzcd}[column sep=large,row sep=large]
X \arrow[d,"\pi"'] \arrow[dr,"\phi_{H-E}"] & \\
\PP^2 \arrow[r,dashed,"\phi_V"'] & \PP^1.
\end{tikzcd}
\]
An explicit model makes the two arrows easy to distinguish:
\[
X\cong
\bigl\{([z_0:z_1:z_2],[u:v])\in\PP^2\times\PP^1:
z_1v=z_2u\bigr\}.
\]
The first projection is \(\pi\), and the second is
\(\phi_{H-E}\). Its fiber over \([u:v]\) is the strict transform of
the line \(z_1v=z_2u\) through \(p\), hence a copy of \(\PP^1\).
The exceptional curve is \(E=\{p\}\times\PP^1\). The map \(\pi\)
contracts it to \(p\), whereas \(\phi_{H-E}|_E\) is an isomorphism
onto the target \(\PP^1\): a point of \(E\) records a tangent
direction at \(p\).

Thus resolving the basepoint makes the map regular, but does not make
it injective. The whole strict transform of each line remains one
fiber. We return to this boundary case in the positivity discussion.
\end{example}
% END CLASSROOM ADDITION

\begin{pauseandthink}
What happens if one starts with a finite subset
\[
A\subseteq M
\]
that is \emph{not} the whole set of lattice points of its convex hull?
Do you expect the orbit closure \(Y_A\) to remember only the convex
hull, or the actual chosen set \(A\)?
\end{pauseandthink}

\subsection{Very ample versus ample}

The map \(\phi_D\) packages all sections of \(\mathcal O(D)\) at once.
There are then three increasingly strong possibilities:

\begin{itemize}
\item The sections may fail to define a morphism everywhere. Then
\(D\) is not basepoint free.

\item The sections may define a morphism
\[
\phi_D:X_\Sigma\to \PP^r,
\]
but different points or tangent directions may still map together.

\item The sections may give a closed embedding. Then \(D\) is very
ample.
\end{itemize}

Ampleness means that this last outcome occurs after replacing \(D\) by
a sufficiently large positive multiple.

So the natural questions become:
\begin{quote}
Can we tell from the support function whether \(\phi_D\) is defined
everywhere? Can we tell when it is an embedding, or becomes one after
scaling?
\end{quote}
Those are precisely the positivity questions answered by convexity and
strict convexity.

\begin{guidingquestion}
In \cref{lem:section-nonzero-fixed-point}, each maximal cone \(\sigma\)
had a distinguished exponent \(m_\sigma\). Why should the condition
\[
m_\sigma\in P_D
\]
mean that the complete linear system has no basepoint at the fixed
point \(x_\sigma\)?

And if two neighboring maximal cones have the same distinguished
exponent, what failure of separation should you expect for the map
\(\phi_D\)?
\end{guidingquestion}
% END INLINED: chapter05-section02-projective-maps-v2.tex

% BEGIN INLINED: chapter05-section03-convexity-positivity-v4.tex
\section{Convexity and positivity}
\label{sec:convexity-positivity}
\index{basepoint-free divisor}
\index{support function!convexity}
\index{support function!strict convexity}
\index{ampleness criterion}

The previous section attached to a basepoint-free divisor a projective
morphism
\[
\phi_D:X_\Sigma\longrightarrow \PP(H^0(X_\Sigma,\mathcal O(D))^\vee).
\]
We now ask how to recognize this situation directly from the fan.
The answer is encoded by the support function. Convexity will tell us
when the map is defined everywhere, and strict convexity will tell us
when the associated projective geometry is positive enough to produce
an embedding after taking multiples.

\begin{agcheatsheet}
Let \(\mathcal L\) be a line bundle on a complete variety \(X\).

A point \(p\in X\) is a \emph{basepoint} of \(\mathcal L\) if every
section in
\[
H^0(X,\mathcal L)
\]
vanishes at \(p\). Equivalently, the evaluation map
\[
H^0(X,\mathcal L)\longrightarrow \mathcal L|_p
\]
is zero.

The line bundle is \emph{generated by global sections}, or
\emph{basepoint free}, if it has no basepoints. This means that near
every point one can find a global section that serves as a local
frame.

It is \emph{very ample} if the complete linear system gives a closed
embedding into projective space. It is \emph{ample} if some positive
tensor power is very ample.
\end{agcheatsheet}

Throughout this section, \(\Sigma\) is a complete fan and
\[
D=\sum_{\rho\in\Sigma(1)}a_\rho D_\rho
\]
is a torus-invariant Cartier divisor. Write
\[
\{m_\sigma\}_{\sigma\in\Sigma(n)}
\]
for its Cartier data and
\[
\psi_D:|\Sigma|\longrightarrow\RR
\]
for its support function.

\subsection{The distinguished section on an affine chart}

Let \(\sigma\in\Sigma(n)\). Since
\[
\langle m_\sigma,u_\rho\rangle=-a_\rho
\qquad
(\rho\in\sigma(1)),
\]
the rational function \(\character{m_\sigma}\) is a local frame of
\(\mathcal O_{X_\Sigma}(D)\) on \(U_\sigma\). Indeed,
\[
\operatorname{div}(\character{m_\sigma})+D
\]
has coefficient zero along every boundary divisor that meets the
closed orbit of \(U_\sigma\).

For any lattice point \(m\in P_D\cap M\), the corresponding section
can be written on this chart as
\[
\character{m}
=
\character{m-m_\sigma}\character{m_\sigma}.
\]
The first factor is regular on \(U_\sigma\), because
\[
m-m_\sigma\in\sigma^\vee.
\]

\begin{lemma}[The section that does not vanish at a fixed point]
\label{lem:section-nonzero-fixed-point}
Let \(x_\sigma=O(\sigma)\) be the torus fixed point corresponding to a
maximal cone \(\sigma\). A monomial section \(\character{m}\), with
\(m\in P_D\cap M\), is nonzero at \(x_\sigma\) if and only if
\[
m=m_\sigma.
\]
Consequently, \(x_\sigma\) is not a basepoint if and only if
\[
m_\sigma\in P_D.
\]
\end{lemma}

\begin{proof}
Relative to the local frame \(\character{m_\sigma}\), the section is
the regular monomial
\[
\character{m-m_\sigma}\in\Bbbk[\sigma^\vee\cap M].
\]
At the distinguished point \(x_\sigma\), a monomial is nonzero only
when its exponent lies in
\[
\sigma^\perp\cap M.
\]
Since \(\sigma\) is full dimensional, this group is zero. Hence the
section is nonzero precisely when
\[
m-m_\sigma=0.
\]
\end{proof}

% BEGIN CLASSROOM ADDITION: notes/20261016.tex (why eigensections matter).
\begin{observation}[The monomial hypothesis matters]
\label{obs:classroom-eigensection-hypothesis}
The preceding fixed-point argument is specific to monomial sections.
A general regular section nonzero at \(x_\sigma\) need only be a
frame near that point: \(1+x\) on \(\AAA^1\) is nonzero at \(0\)
but vanishes at \(-1\). Nor do boundary orders alone determine the
value of an arbitrary section at a fixed point: \(x+y\) vanishes at
the origin of \(\AAA^2\), although it has order zero along both
coordinate axes. Passing to a nonvanishing weight component is what
makes the monomial criterion applicable to global generation.
\end{observation}
% END CLASSROOM ADDITION

This lemma gives a useful interpretation of the Cartier datum:
\[
\boxed{
 m_\sigma
 \text{ is the exponent of the monomial section that should
 trivialize }\mathcal O(D)\text{ on }U_\sigma.
}
\]
The only question is whether that exponent actually lies in \(P_D\),
so that the local frame extends to a global section.

% BEGIN CLASSROOM ADDITION: notes/20261016.tex (the local model x^2/y).
\begin{example}[A local equation and its reciprocal frame]
\label{ex:classroom-reciprocal-frame}
Suppose \(U_\sigma\simeq\AAA^2\), with boundary divisors
\(V(x),V(y)\), and
\[
D|_{U_\sigma}=2V(x)-V(y)
=\operatorname{div}\!\left(\frac{x^2}{y}\right).
\]
The local defining equation has a zero of order two along \(V(x)\)
and a pole of order one along \(V(y)\). Its reciprocal is a frame of
\(\mathcal O(D)\): in the rational-function realization used here,
\[
m_\sigma=(-2,1),\qquad e_\sigma=\character{m_\sigma}=\frac{y}{x^2}.
\]
Every regular section on this chart is
\[
f=p(x,y)e_\sigma=p(x,y)\frac{y}{x^2},
\qquad p\in\Bbbk[x,y].
\]
Its value in the line-bundle fiber at \(x_\sigma=(0,0)\) is nonzero
exactly when \(p(0,0)\ne0\). In that case,
\[
\ord_{V(x)}(f)=-2,\qquad \ord_{V(y)}(f)=1.
\]
The rational function \(f\) can therefore have a pole while still
representing a regular section of \(\mathcal O(D)\). The coefficient
\(p\), not the numerical value of \(f\), is evaluated in this frame.
The monomial frame extends globally exactly when
\((-2,1)\in P_D\); rays outside \(\sigma\) supply the remaining
inequalities for that extension.
\end{example}
% END CLASSROOM ADDITION

\subsection{Convex support functions and global generation}

\begin{definition}[Convexity in the toric convention]
\label{def:toric-convex-support-function}
The support function \(\psi_D\) is \emph{convex} if, for every maximal
cone \(\sigma\),
\[
\psi_D(u)\leq\langle m_\sigma,u\rangle
\qquad
\text{for every }u\in N_\RR.
\]
Equivalently,
\[
\psi_D(u)
=
\min_{\sigma\in\Sigma(n)}
\langle m_\sigma,u\rangle.
\]
\end{definition}

Each linear piece agrees with \(\psi_D\) on its own cone. Convexity
says that when this piece is extended across all of \(N_\RR\), the
graph of \(\psi_D\) never rises above it.

\begin{warningbox}
Because \(\psi_D\) is written as a \emph{minimum} of linear
functions, some books on convex analysis would call it concave or
upper convex. We follow the standard toric convention used by
Cox--Little--Schenck: this condition is called convex.
\end{warningbox}

% BEGIN CLASSROOM ADDITION: notes/20261016.tex (the classroom sign convention).
\begin{observation}[Translating the classroom maximum formula]
\label{obs:classroom-support-sign}
If the local defining equation is written as
\(\character{q_\sigma}\), then \(q_\sigma=-m_\sigma\).
The opposite support function \(\eta_D=-\psi_D\) has
\(\eta_D(u_\rho)=a_\rho\), and
\[
\psi_D(u)=\min_\sigma\langle m_\sigma,u\rangle
\quad\Longleftrightarrow\quad
\eta_D(u)=\max_\sigma\langle q_\sigma,u\rangle.
\]
This is the maximum formula used in class, with the local-equation
exponents in place of the frame exponents. Throughout this chapter,
\(\psi_D\) retains the minimum convention. Both formulas express the
same global-generation condition.
\end{observation}
% END CLASSROOM ADDITION

\begin{theorem}[Basepoint freeness equals convexity]
\label{thm:basepoint-free-convex}
The following conditions are equivalent:

\begin{enumerate}[label=(\alph*),wide=0pt]
\item \(D\) is basepoint free.

\item
\[
m_\sigma\in P_D
\qquad
\text{for every }\sigma\in\Sigma(n).
\]

\item
\[
\psi_D(u)
=
\min_{\sigma\in\Sigma(n)}
\langle m_\sigma,u\rangle
\qquad
\text{for every }u\in N_\RR.
\]

\item The support function \(\psi_D\) is convex.
\end{enumerate}
\end{theorem}

\noindent
\textit{Reference.}
This is Cox--Little--Schenck, Theorem~6.1.7.

\begin{proof}
Suppose first that every \(m_\sigma\) lies in \(P_D\). By
\cref{thm:lattice-points-global-sections},
\(\character{m_\sigma}\) is a global section. It is a local frame on
\(U_\sigma\), so these sections generate \(\mathcal O(D)\) on the
affine cover
\[
X_\Sigma=\bigcup_{\sigma\in\Sigma(n)}U_\sigma.
\]
Thus \textup{(b)} implies \textup{(a)}.

Conversely, suppose \(D\) is basepoint free. At the fixed point
\(x_\sigma\), some global section is nonzero. Decomposing that section
into its monomial weight components, at least one monomial section is
nonzero there. By \cref{lem:section-nonzero-fixed-point}, its exponent
must be \(m_\sigma\). Hence \(m_\sigma\in P_D\), proving
\textup{(a)} implies \textup{(b)}.

Now fix maximal cones \(\sigma,\tau\), and let \(u\in\tau\). If
\(m_\sigma\in P_D\), then the inequalities defining \(P_D\) imply
\[
\langle m_\sigma,u\rangle
\geq
\langle m_\tau,u\rangle
=
\psi_D(u).
\]
Since equality holds when \(u\in\sigma\), this is equivalent to
\[
\psi_D(u)=\min_\sigma\langle m_\sigma,u\rangle.
\]
Thus \textup{(b)}, \textup{(c)}, and \textup{(d)} are equivalent.
\end{proof}

The theorem turns a local geometric question into the shape of one
piecewise-linear function:
\[
\boxed{
\mathcal O(D)\text{ is globally generated}
\quad\Longleftrightarrow\quad
\psi_D\text{ bends convexly across every wall.}
}
\]

% BEGIN CLASSROOM ADDITION: notes/20261016.tex (the three linear pieces on P^2).
\begin{example}[The support function of a line on \(\PP^2\)]
\label{ex:classroom-p2-support-function}
For the ray generators \(u_0=-e_1-e_2,u_1=e_1,u_2=e_2\), take
\(D=D_0\). Its polytope is \(\Delta_2\), and its Cartier data and
linear pieces are
\[
\begin{array}{c|c|c}
\sigma&m_\sigma&\psi_D(x,y)\text{ on }\sigma\\ \hline
\Cone{e_1,e_2}&(0,0)&0\\
\Cone{e_2,-e_1-e_2}&(1,0)&x\\
\Cone{-e_1-e_2,e_1}&(0,1)&y
\end{array}
\]
All three Cartier data lie in \(\Delta_2\). On each cone, the
listed linear piece is the smallest of \(0,x,y\), so
\[
\psi_D(x,y)=\min(0,x,y).
\]
The classroom expression \(\max(0,-x,-y)\) is its negative, as in
\cref{obs:classroom-support-sign}. Along a common ray the neighboring
formulas agree, making continuity and the convexity criterion
visible in this example.
\end{example}
% END CLASSROOM ADDITION

\begin{observation}
There is also a numerical name for this condition. A Cartier divisor
is \emph{nef} if
\[
D\cdot C\geq0
\]
for every complete curve \(C\). On a toric variety, a Cartier divisor
is nef exactly when it is basepoint free; it is enough to test the
torus-invariant complete curves.

This is Cox--Little--Schenck, Theorem~6.3.12. No intersection theory
is needed for the arguments in this section, so this terminology may
be viewed as an optional preview.
\end{observation}

\subsection{Strict convexity and ampleness}

Convexity allows two adjacent maximal cones to carry the same linear
piece. Geometrically, the support function then fails to notice the
wall between them, and the projective map from the previous section
cannot separate the geometry on the two sides. Ampleness requires a
genuine bend at every wall.

\begin{definition}[Strict convexity]
\label{def:strictly-convex-support-function}
A convex support function \(\psi_D\) is \emph{strictly convex} if, for
every maximal cone \(\sigma\),
\[
\psi_D(u)=\langle m_\sigma,u\rangle
\quad\Longleftrightarrow\quad
 u\in\sigma.
\]
Equivalently, neighboring maximal cones have different Cartier data:
\[
m_\sigma\neq m_{\sigma'}
\]
whenever \(\sigma\) and \(\sigma'\) meet along a wall.
\end{definition}

\begin{figure}[ht]
\centering
\begin{tikzpicture}[scale=0.95]
  % Convex but not strict
  \begin{scope}[xshift=-4.4cm]
    \draw[->,black!60] (-2.85,0)--(2.85,0) node[below right] {slice};
    \draw[->,black!60] (0,-1.55)--(0,2.15)
      node[right] {\(\psi\)};

    \foreach \x in {-1.2,0,1.2} {
      \draw[black!35,dashed] (\x,-1.25)--(\x,1.85);
      \node[black!55,below] at (\x,-1.25) {wall};
    }

    \draw[very thick,DictionaryBlue]
      (-2.5,-0.9)--(-1.2,0.55)--(1.2,0.55)--(2.5,-0.9);

    \node[fill=DictionaryBlue!10,rounded corners,inner sep=2pt]
      at (0,1.25) {one linear piece crosses the middle wall};
    \node[font=\bfseries] at (0,2.55)
      {convex, not strict};
    \node[align=center] at (0,-2.05)
      {the map is defined,\\but one wall is invisible};
  \end{scope}

  % Strictly convex
  \begin{scope}[xshift=4.4cm]
    \draw[->,black!60] (-2.85,0)--(2.85,0) node[below right] {slice};
    \draw[->,black!60] (0,-1.55)--(0,2.15)
      node[right] {\(\psi\)};

    \foreach \x in {-1.2,0,1.2} {
      \draw[black!35,dashed] (\x,-1.25)--(\x,1.85);
      \node[black!55,below] at (\x,-1.25) {wall};
    }

    \draw[very thick,PredictionTeal]
      (-2.5,-0.95)--(-1.2,0.35)--(0,1.1)--(1.2,0.35)--(2.5,-0.95);

    \node[fill=PredictionTeal!10,rounded corners,inner sep=2pt]
      at (0,1.65) {a bend occurs at every wall};
    \node[font=\bfseries] at (0,2.55)
      {strictly convex};
    \node[align=center] at (0,-2.05)
      {each wall changes the slope,\\so no wall is forgotten};
  \end{scope}
\end{tikzpicture}
\caption{A transverse slice through the graph of a support function.
Both panels use the same three walls. On the left, the middle wall is
invisible to the support function because one linear piece continues
straight across it. On the right, the slope changes at every wall.}
\label{fig:convex-versus-strict-support}
\end{figure}

\begin{theorem}[Ampleness equals strict convexity]
\label{thm:ample-strictly-convex}
Let \(D\) be a torus-invariant Cartier divisor on the complete toric
variety \(X_\Sigma\). Then
\[
\boxed{
D\text{ is ample}
\quad\Longleftrightarrow\quad
\psi_D\text{ is strictly convex}.
}
\]
\end{theorem}

\noindent
\textit{Reference.}
This is Cox--Little--Schenck, Theorem~6.1.14.

\begin{proof}[Conceptual proof sketch]
If a positive multiple of \(D\) gives an embedding, then the resulting
projective map must separate the torus-invariant curve attached to
any wall. If two neighboring maximal cones had the same Cartier
datum, the corresponding support function would fail to bend at that
wall, and the linear system would not detect that invariant curve.
Thus ampleness forces a genuine change of slope everywhere.

Conversely, strict convexity ensures that different maximal cones have
different Cartier data, so the polytope has enough distinct vertices
to encode the fan faithfully. The projective map attached to a large
multiple of \(D\) then separates the toric strata, and one obtains a
closed embedding. The complete proof is given in the cited theorem.
\end{proof}

\begin{corollary}[Projectivity criterion]
\label{cor:projective-strict-support-function}
A complete toric variety \(X_\Sigma\) is projective if and only if its
fan admits an integral strictly convex support function.
\end{corollary}

This criterion is easy to state but powerful: completeness is visible
from whether the cones cover \(N_\RR\), while projectivity requires a
compatible strictly convex shape on those cones.

\subsection{Surface laboratory: the positivity interval on a blowup}

Let
\[
X=\Blp{\PP^2}{p}
\]
and consider
\[
D=dD_0-rE
\]
with \(d>0\) and \(0\leq r\leq d\). As in
\cref{fig:blowup-truncated-section-polytope},
\[
P_D
=
\left\{
(a,b)\in\RR^2:
 a,b\geq0,\quad r\leq a+b\leq d
\right\}.
\]

For the four maximal cones, the Cartier data are
\[
(0,d),
\qquad
(0,r),
\qquad
(r,0),
\qquad
(d,0).
\]
Each lies in \(P_D\), so \(D\) is basepoint free throughout the range
\[
0\leq r\leq d.
\]
They are distinct exactly when
\[
0<r<d.
\]
Therefore
\[
\boxed{
 dH-rE\text{ is basepoint free for }0\leq r\leq d,
 \qquad
 dH-rE\text{ is ample for }0<r<d.
}
\]

\begin{figure}[ht]
\centering
\begin{tikzpicture}[scale=0.82]
  % 2H
  \begin{scope}[xshift=-5.2cm]
    \draw[->,black!60] (-0.35,0)--(2.65,0) node[below] {\(a\)};
    \draw[->,black!60] (0,-0.35)--(0,2.65) node[left] {\(b\)};
    \filldraw[fill=DictionaryBlue!13,draw=DictionaryBlue,very thick]
      (0,0)--(2,0)--(0,2)--cycle;
    \draw[WarningRed,dashed,very thick] (-0.3,0.3)--(0.3,-0.3);
    \fill (0,0) circle (2.3pt);
    \fill (2,0) circle (2.3pt);
    \fill (0,2) circle (2.3pt);
    \node[WarningRed,fill=white,inner sep=1pt,anchor=east] at (-0.08,0.78)
      {redundant wall};
    \draw[WarningRed,->,thick] (-0.14,0.67)--(-0.02,0.09);
    \node[font=\bfseries] at (0.95,3.05) {\(2H\)};
    \node[align=center,font=\small] at (0.95,-0.68)
      {convex, not strict};
  \end{scope}

  % 2H-E
  \begin{scope}
    \draw[->,black!60] (-0.2,0)--(2.65,0) node[below] {\(a\)};
    \draw[->,black!60] (0,-0.2)--(0,2.65) node[left] {\(b\)};
    \filldraw[fill=PredictionTeal!14,draw=PredictionTeal,very thick]
      (1,0)--(2,0)--(0,2)--(0,1)--cycle;
    \foreach \p in {(1,0),(2,0),(0,2),(0,1)} {
      \fill \p circle (2.3pt);
    }
    \draw[PredictionTeal,very thick]
      (1,0)--(0,1)
      node[midway,above,sloped,fill=white,inner sep=1pt,text=PredictionTeal,yshift=5pt]
      {new active wall};
    \node[font=\bfseries] at (0.95,3.05) {\(2H-E\)};
    \node[align=center,font=\small] at (0.95,-0.68)
      {strictly convex};
  \end{scope}

  % 2H-2E
  \begin{scope}[xshift=5.2cm]
    \draw[->,black!60] (-0.2,0)--(2.65,0) node[below] {\(a\)};
    \draw[->,black!60] (0,-0.2)--(0,2.65) node[left] {\(b\)};
    \draw[SurfaceGreen,very thick,line width=2.5pt]
      (2,0)--(0,2)
      node[midway,above,sloped,fill=white,inner sep=1pt,text=SurfaceGreen,yshift=5pt]
      {\(a+b=2\)};
    \fill (2,0) circle (2.3pt);
    \fill (0,2) circle (2.3pt);
    \node[font=\bfseries] at (0.95,3.05) {\(2H-2E\)};
    \node[align=center,font=\small] at (0.95,-0.68)
      {convex, not strict};
  \end{scope}
\end{tikzpicture}
\caption{Positivity for divisors \(dH-rE\) on the blowup. For
\(2H\), the exceptional wall is redundant. For \(2H-E\), the new wall
is active and all four Cartier data are distinct. For \(2H-2E\), the
polytope collapses to a segment.}
\label{fig:blowup-convexity-positivity}
\end{figure}

The two boundary cases fail strict convexity for different visible
reasons:

\begin{itemize}
\item When \(r=0\), the inequality from the exceptional ray is
redundant. The support function does not bend along the newly inserted
wall.

\item When \(r=d\), the polytope loses dimension. The linear system
remembers only the pencil of lines through the blown-up point.
\end{itemize}

% BEGIN CLASSROOM ADDITION: notes/20261014.tex (the two boundary maps).
The classroom projection in \cref{ex:classroom-projection-pencil}
also makes both boundary cases explicit. For every \(k\geq1\),
\[
\phi_{kH}=\nu_k\circ\pi,\qquad
\phi_{k(H-E)}=\nu_k\circ\phi_{H-E},
\]
where the Veronese maps are those of \(\PP^2\) and \(\PP^1\),
respectively. Thus every positive multiple of \(H\) still contracts
\(E\), and every positive multiple of \(H-E\) still contracts the
strict transforms of the lines through \(p\). These divisors are
basepoint free without becoming very ample after scaling.
% END CLASSROOM ADDITION

\begin{pauseandthink}
On \(\PP^1\times\PP^1\), show that the polytope of
\(aF_1+bF_2\) is a rectangle with side lengths \(a\) and \(b\).
Use its shape to determine when the divisor is basepoint free and when
it is ample. What collapses when one coefficient is zero?
\end{pauseandthink}

For a basepoint-free divisor, the Cartier data lie in \(P_D\). When
different maximal cones have the same Cartier datum, the polytope
forgets the wall between them. Strict convexity is exactly the
condition that no wall is forgotten.

\begin{threedictionaries}
At this point, the three descriptions of positivity line up.

\medskip

\begin{center}
\small
\begin{tabularx}{\linewidth}{@{}
  >{\raggedright\arraybackslash}X
  >{\raggedright\arraybackslash}X
  >{\raggedright\arraybackslash}X@{}}
\toprule
\textbf{Algebra}
&
\textbf{Geometry}
&
\textbf{Combinatorics}
\\
\midrule
\(\character{m_\sigma}\) extends to a global section
&
\(x_\sigma\) is not a basepoint
&
\(m_\sigma\in P_D\cap M\)
\\[1mm]
global sections generate \(\mathcal O(D)\)
&
\(\phi_D\) is defined everywhere
&
\(\psi_D\) is convex
\\[1mm]
sections of \(kD\) separate toric strata for \(k\gg0\)
&
\(D\) is ample
&
\(\psi_D\) is strictly convex
\\
\bottomrule
\end{tabularx}
\end{center}

\[
\boxed{
\text{convexity makes the map exist;\quad
strict convexity makes it eventually embed.}
}
\]
\end{threedictionaries}

\begin{guidingquestion}
Given a polytope \(P\subseteq M_\RR\), how can one construct a fan in
\(N_\RR\) from the directions in which each face of \(P\) is
minimized?

For a basepoint-free divisor \(D\), how does that fan compare with the
original fan \(\Sigma\), and why does strict convexity force the two
fans to agree?
\end{guidingquestion}
% END INLINED: chapter05-section03-convexity-positivity-v4.tex

% BEGIN INLINED: chapter05-section04-complete.tex
\section{Polytopes and normal fans}
\label{sec:polytopes-normal-fans}
\index{normal cone}
\index{normal fan}
\index{polytope!faces}
\index{contraction!from a normal fan}

A polytope does more than provide lattice points for projective
coordinates. It also remembers a fan.

The guiding idea is simple. Choose a direction
\[
u\in N_\RR.
\]
The linear function
\[
m\longmapsto\langle m,u\rangle
\]
reaches its minimum on a face of the polytope. Directions that select
the same face form a cone. As the face changes, these cones fit
together into the \emph{normal fan}.

\subsection{Which face does a direction see?}

Let \(P\subseteq M_\RR\) be a full-dimensional lattice polytope. For
\(u\in N_\RR\), define the minimizing face
\[
\operatorname{face}_u(P)
=
\left\{
 m\in P:
 \langle m,u\rangle
 =
 \min_{p\in P}\langle p,u\rangle
\right\}.
\]

The vector \(u\) should be thought of as an inward-pointing normal:
it pushes a supporting hyperplane toward \(P\) until that hyperplane
first meets the polytope.

\begin{definition}[Normal cone]
\label{def:normal-cone-face}
Let \(Q\preceq P\) be a face. The \emph{normal cone} of \(P\) along
\(Q\) is
\[
\sigma_Q
=
\left\{
 u\in N_\RR:
 Q\subseteq\operatorname{face}_u(P)
\right\}.
\]
Equivalently, \(u\in\sigma_Q\) precisely when
\[
\langle q,u\rangle
\leq
\langle p,u\rangle
\qquad
\text{for every }q\in Q,
p\in P.
\]
\end{definition}

The relative interior has the sharper interpretation
\[
u\in\relint(\sigma_Q)
\quad\Longleftrightarrow\quad
\operatorname{face}_u(P)=Q.
\]
Thus the closed cone allows a larger minimizing face, while its
relative interior remembers exactly \(Q\).

Suppose the facets of \(P\) have equations
\[
P
=
\left\{
 m\in M_\RR:
 \langle m,u_F\rangle\geq-a_F
 \text{ for every facet }F
\right\},
\]
where \(u_F\in N\) is the primitive inward normal to \(F\). Then
\[
\boxed{
\sigma_Q
=
\Cone{u_F:Q\subseteq F}.
}
\]
In particular:
\[
\begin{array}{c@{\qquad\longleftrightarrow\qquad}c}
\text{facet }F
&
\text{ray }\Cone{u_F},
\\[1mm]
\text{vertex }v
&
\text{maximal cone generated by the normals at }v.
\end{array}
\]

\subsection{The normal fan}

\begin{theorem}[Faces assemble into a fan]
\label{thm:normal-fan-polytope}
Let \(P\subseteq M_\RR\) be a full-dimensional lattice polytope. Then
\[
\Sigma_P
=
\{\sigma_Q:Q\preceq P\}
\]
is a complete fan in \(N_\RR\), called the \emph{normal fan} of
\(P\).
\end{theorem}

\noindent
\textit{Reference.}
This is Cox--Little--Schenck, Theorem~2.3.2; see also
Proposition~6.2.3.

\begin{proof}[Conceptual proof sketch]
If \(Q\subseteq Q'\), then every supporting direction that minimizes
on \(Q'\) also minimizes on \(Q\). Hence inclusions reverse:
\[
Q\subseteq Q'
\quad\Longrightarrow\quad
\sigma_{Q'}\preceq\sigma_Q.
\]
The faces of a normal cone therefore come from larger faces of the
polytope.

For two faces \(Q,Q'\), the intersection
\[
\sigma_Q\cap\sigma_{Q'}
\]
consists of directions minimizing on both. This is the normal cone of
the smallest face containing \(Q\) and \(Q'\), so the intersection is
a face of each cone.

Finally, every \(u\in N_\RR\) minimizes somewhere on the compact
polytope \(P\). Therefore \(u\) lies in the normal cone of its
minimizing face, proving completeness.
\end{proof}

The face--cone correspondence reverses dimensions and inclusions:
\[
\boxed{
\dim Q+\dim\sigma_Q=\dim P.
}
\]
For a polygon, this gives the compact dictionary
\[
\begin{array}{c|c}
\text{face of }P
&
\text{cone of }\Sigma_P
\\
\hline
\text{vertex}
&
\text{two-dimensional cone}
\\
\text{edge}
&
\text{ray}
\\
P
&
\{0\}.
\end{array}
\]

\begin{figure}[!htb]
\centering
\begin{tikzpicture}[scale=0.78]
  % Polytope
  \begin{scope}[xshift=-4.8cm]
    \draw[->,black!55] (-0.3,0)--(3.1,0) node[below] {\(a\)};
    \draw[->,black!55] (0,-0.3)--(0,3.1) node[left] {\(b\)};

    \filldraw[fill=PredictionTeal!10,draw=black!75,very thick]
      (1,0)--(2,0)--(0,2)--(0,1)--cycle;

    \draw[very thick,SurfaceGreen]
      (1,0)--(2,0)
      node[midway,above=3pt,text=SurfaceGreen] {\(e_2\)};
    \draw[very thick,ExampleGold]
      (2,0)--(0,2)
      node[midway,above=4pt,sloped,text=ExampleGold] {\(-e_1-e_2\)};
    \draw[very thick,DictionaryBlue]
      (0,2)--(0,1)
      node[midway,left=5pt,text=DictionaryBlue] {\(e_1\)};
    \draw[very thick,PredictionTeal]
      (0,1)--(1,0)
      node[midway,below=4pt,sloped,text=PredictionTeal] {\(e_1+e_2\)};

    \foreach \p in {(1,0),(2,0),(0,2),(0,1)} {
      \fill \p circle (2.2pt);
    }

    \node[font=\bfseries] at (1.1,3.55)
      {the polytope \(P_{2H-E}\)};
    \node[align=center,font=\small] at (1.1,-1.05)
      {each edge contributes\\its inward normal};
  \end{scope}

  % Normal fan
  \begin{scope}[xshift=4.6cm]
    \fill[DictionaryBlue!9]
      (0,0)--(2.6,0)--(1.75,1.75)--cycle;
    \fill[PredictionTeal!9]
      (0,0)--(1.75,1.75)--(0,2.6)--cycle;
    \fill[SurfaceGreen!9]
      (0,0)--(0,2.6)--(-1.85,-1.85)--cycle;
    \fill[ExampleGold!10]
      (0,0)--(-1.85,-1.85)--(2.6,0)--cycle;

    \draw[very thick,DictionaryBlue,-{Latex[length=2.2mm]}]
      (0,0)--(2.95,0);
    \draw[very thick,PredictionTeal,-{Latex[length=2.2mm]}]
      (0,0)--(2.2,2.2);
    \draw[very thick,SurfaceGreen,-{Latex[length=2.2mm]}]
      (0,0)--(0,2.95);
    \draw[very thick,ExampleGold,-{Latex[length=2.2mm]}]
      (0,0)--(-2.15,-2.15);

    \node[DictionaryBlue,below] at (2.7,0) {\(e_1\)};
    \node[PredictionTeal,above right] at (1.95,1.95)
      {\(e_1+e_2\)};
    \node[SurfaceGreen,left] at (0,2.7) {\(e_2\)};
    \node[ExampleGold,above left] at (-1.95,-1.95)
      {\(-e_1-e_2\)};

    \node[font=\bfseries] at (0.35,3.55)
      {the normal fan \(\Sigma_P\)};
    \node[align=center,font=\small] at (0.35,-2.75)
      {the normals are reassembled\\at the origin};
  \end{scope}
\end{tikzpicture}
\caption{The truncated triangle for \(2H-E\) and its normal fan. Each
edge contributes its primitive inward normal. At a vertex, the normals
of the two incident edges generate the corresponding maximal cone.}
\label{fig:truncated-triangle-normal-fan}
\end{figure}

\begin{prediction}
Before looking at the assembled fan, focus on the vertex \((1,0)\).
Which two inward facet normals should generate its normal cone?

Choose a vector in the interior of your predicted cone and determine
where its linear function is minimized on \(P\). Then compare your
answer with \cref{fig:truncated-triangle-normal-fan}.
\end{prediction}

% BEGIN CLASSROOM ADDITION: notes/20261016.tex (the trapezoid vertex charts).
\subsection{Recovering a variety from its vertex charts}

\begin{example}[The five-monomial trapezoid]
\label{ex:classroom-trapezoid-charts}
The classroom question ``which variety does a polytope produce?'' can
be answered directly in affine coordinates. Consider
\[
P=\Conv\{(0,0),(2,0),(1,1),(0,1)\}
\]
and its five lattice points, ordered as
\[
m_0=(0,0),\quad m_1=(1,0),\quad m_2=(0,1),\quad
m_3=(2,0),\quad m_4=(1,1).
\]
Writing \((s,t)\) for coordinates on the torus, set
\[
\phi_P(s,t)=[1:s:t:s^2:st],\qquad
Y=\overline{\phi_P(T)}\subseteq\PP^4.
\]
Let \(Y_i=Y\cap\{x_i\ne0\}\). On the chart \(i=0\),
the image has coordinates \((s,t,s^2,st)\). Its closure is the graph
of the polynomial map \((s,t)\mapsto(s^2,st)\), so
\[
Y_0=V(z_3-z_1^2,z_4-z_1z_2)\cong\AAA^2,
\qquad z_j=x_j/x_0.
\]
The remaining vertex charts work in the same way. In each tuple below,
the coordinates \(x_j/x_i\) appear in increasing order of \(j\ne i\).
\[
\begin{array}{c|c|c|c}
i&m_i&\text{parameters }(u,v)&\text{normalized image}\\ \hline
0&(0,0)&(s,t)&(u,v,u^2,uv)\\
2&(0,1)&(s,t^{-1})&(v,uv,u^2v,u)\\
3&(2,0)&(s^{-1},t/s)&(u^2,u,uv,v)\\
4&(1,1)&(s^{-1},s/t)&(u^2v,uv,u,v)
\end{array}
\]
Each tuple contains \(u,v\) themselves, with the other entries
polynomial in them. Thus each of these four charts is an affine plane.

Why do normal cones appear? Dividing by the coordinate of a vertex
\(m_i\) replaces each character \(\character{m}\) by
\(\character{m-m_i}\). Their exponents generate the tangent cone
\[
\sigma_i^\vee=\Cone{P-m_i}\subseteq M_\RR.
\]
Its dual is precisely the inward normal cone at \(m_i\):
\[
\sigma_i=\{u\in N_\RR:
\langle m-m_i,u\rangle\geq0\text{ for every }m\in P\}.
\]
The tangent cone lies in the character space, whereas the normal
cone lies in the dual space. In the present example the normal cones
are
\[
\begin{aligned}
\sigma_0&=\Cone{e_1,e_2},&
\sigma_2&=\Cone{e_1,-e_2},\\
\sigma_3&=\Cone{e_2,-e_1-e_2},&
\sigma_4&=\Cone{-e_1-e_2,-e_2}.
\end{aligned}
\]
At each vertex the two parameter exponents form a lattice basis, so
the displayed coordinate ring is exactly
\(\Bbbk[\sigma_i^\vee\cap M]\).

The lattice point \((1,0)\) is not a vertex. Its chart has coordinates
\[
(s^{-1},t/s,s,t)=(u^{-1},v,u,uv),
\]
and therefore
\[
Y_1\cong\Spec\Bbbk[u,u^{-1},v]\cong\Gm\times\AAA^1.
\]
The relation \(x_1^2=x_0x_3\) on \(Y\) gives
\(Y_1=Y_0\cap Y_3\). Consequently the four vertex charts cover \(Y\).
They glue through their common torus characters; for example, on
\(Y_0\cap Y_3\) the coordinate change is
\((s,t)\mapsto(s^{-1},t/s)\). These are the affine charts and gluing
maps of the normal fan, so
\[
Y\cong X_{\Sigma_P}.
\]

This is also the earlier \(2H-E\) example in different lattice
coordinates: the affine lattice transformation
\[
(r,q)\longmapsto(a,b)=(2-r-q,r)
\]
sends \(P\) to
\(\{a,b\geq0:1\leq a+b\leq2\}=P_{2H-E}\).
The linear part of this transformation has determinant one. The chart
calculation therefore explains concretely why the five sections of
\(2H-E\) embed the blowup into \(\PP^4\).
\end{example}

\begin{observation}[The semigroup check]
\label{obs:classroom-vertex-semigroups}
For a finite set \(A\) affinely spanning \(M_\RR\) and a vertex
\(v\) of its convex hull, the
corresponding affine chart of the orbit closure has coordinate ring
\[
\Bbbk[S_v],\qquad S_v=\NN\{m-v:m\in A\}.
\]
By \cref{thm:normality-saturation-cone}, its cone alone need not determine this ring: the semigroup may have
missing lattice points. When the differences in \(A\) generate \(M\),
saturating these semigroups gives the affine charts of
\(X_{\Sigma_{\Conv(A)}}\), hence the normalization of
the orbit closure. If they generate a proper sublattice, that
sublattice is the character lattice of the image torus and must be
used in the normalization statement. In the trapezoid example the
parameter calculations verify saturation in \(M\) at every vertex,
so no normalization is needed.
\end{observation}
% END CLASSROOM ADDITION

% BEGIN CLASSROOM ADDITION: notes/20261019.tex (local functions and very ampleness).
\begin{observation}[Local functions versus normalized global sections]
In the trapezoid example, let \(D\) be the invariant Cartier divisor
with \(P_D=P\), and take
\(\sigma=\Cone{e_2,-e_1-e_2}\), whose Cartier datum is
\(m_\sigma=(2,0)\). The monomial frame is
\label{obs:classroom-oct19-local-global}
\(e_\sigma=s^2s_D\), and division by it identifies
\[
H^0(U_\sigma,\mathcal O(D))
\xrightarrow{\ \widetilde s\mapsto\widetilde s/e_\sigma\ }
\Bbbk[\sigma^\vee\cap M]
=\Bbbk[s^{-1},t/s].
\]
The translated lattice points \((P_D\cap M)-m_\sigma\) record the
normalized \emph{global} sections; the whole cone records all local
regular monomials. For example, \(s^{-3}\) is regular on this chart,
but
\[
s^{-3}e_\sigma=s^{-1}s_D
\]
does not extend to a global regular section, since \((-1,0)\notin P_D\).
Thus trivializing the bundle and extending a section globally are
different operations.
\end{observation}

\subsection{When the vertex coordinates give an embedding}

The semigroup check also makes the distinction between ample and very
ample precise. Let \(P\subseteq M_\RR\) be a full-dimensional lattice
polytope, and let
\(D_P\) be its ample Cartier divisor on \(X_{\Sigma_P}\). At a vertex
\(v\), write
\[
S_v=\NN\{m-v:m\in P\cap M\},\qquad
\sigma_v^\vee=\Cone{P-v}.
\]
The normalized projective coordinates generate \(\Bbbk[S_v]\), whereas
the coordinate ring of the prescribed normal affine chart is
\(\Bbbk[\sigma_v^\vee\cap M]\). Consequently,
\[
\boxed{
D_P\text{ is very ample}
\quad\Longleftrightarrow\quad
S_v=\sigma_v^\vee\cap M\text{ at every vertex }v.
}
\]
Indeed, the nonvanishing locus of the vertex section is
\(U_{\sigma_v}\). On this open set, a closed embedding is equivalent
to surjectivity of the coordinate-ring map, which is exactly the
displayed semigroup equality. These vertex opens cover the variety
and the corresponding affine target opens cover its image. That
image is closed because \(X_{\Sigma_P}\) is projective.
Equality of the \emph{real cones} alone does not ensure equality of
their integral semigroups.

Every lattice polygon satisfies this condition: it admits a
triangulation into unimodular lattice triangles, and dilating those
triangles shows that every lattice point of \(kP\) is a sum of \(k\)
lattice points of \(P\). Applied near a vertex, this also gives the
displayed tangent-semigroup equality. In dimension three the equality
can fail, even for an ample divisor.

\begin{example}[An ample divisor that needs to be doubled]
\label{ex:classroom-oct19-reeve}
For this example work over \(\CC\), as in the lecture, and take
\[
M=\ZZ^3,\qquad w=(1,1,2),\qquad
P=\Conv\{0,e_1,e_2,w\}.
\]
This is the height-two Reeve tetrahedron. Its only lattice points are
its four vertices, so its complete linear system restricts on the
torus to
\[
\phi_{D_P}(s,t,u)=[1:s:t:stu^2].
\]
The points \((s,t,u)\) and \((s,t,-u)\) have the same image. Moreover,
every \([1:a:b:c]\) with \(abc\ne0\) is obtained by choosing
\(u^2=c/(ab)\). Hence the projective image is all of \(\PP^3\).
The exponent differences generate
\[
M'=\ZZ e_1+\ZZ e_2+\ZZ w
=\{(a,b,c)\in M:c\text{ is even}\},
\]
an index-two sublattice. Thus the map
\(X_{\Sigma_P}\to\PP^3\) is finite of degree two, and is not an
embedding. This is the lattice issue in
Observation~\ref{obs:classroom-vertex-semigroups}, rather than a normalization of
\(\PP^3\).

At the vertex \(0\), the missing lattice point
\[
h=(1,1,1)=\tfrac12(e_1+e_2+w)
\]
lies in the tangent cone but not in \(S_0\). Its character is \(stu\).
The cone semigroup is generated by \(e_1,e_2,w,h\), so the normal-fan
chart is
\[
\Spec\CC[s,t,stu^2,stu]
\cong\Spec\CC[a,b,c,d]/(d^2-abc).
\]
Its origin is singular, whereas the image chart
\(\Spec\CC[s,t,stu^2]\) is an affine three-space. In particular,
smoothness of the image does not imply smoothness of the prescribed
normal-fan variety.

% Classroom addition from notes/20261019.tex, Diagrams 5--6.
\begin{figure}[htbp]
\centering
\begin{tikzpicture}[font=\small,line join=round]
  % Equal display sizes; the coordinates in the right panel are twice as large.
  \begin{scope}[
    x={(1.7325cm,-0.42cm)},
    y={(-0.9975cm,0.315cm)},
    z={(0.105cm,1.155cm)}]
    \coordinate (O) at (0,0,0);
    \coordinate (A) at (1,0,0);
    \coordinate (B) at (0,1,0);
    \coordinate (W) at (1,1,2);
    \fill[DictionaryBlue!7] (O)--(A)--(W)--cycle;
    \fill[DictionaryBlue!13] (O)--(B)--(W)--cycle;
    \draw[DictionaryBlue,thick,dashed] (A)--(B);
    \draw[DictionaryBlue,very thick] (O)--(A)--(W)--(B)--cycle;
    \draw[DictionaryBlue,thick] (O)--(W);
    \foreach \point in {O,A,B,W}
      \fill[DictionaryBlue] (\point) circle (2pt);
    \node[below left] at (O) {\(0\)};
    \node[below right] at (A) {\(e_1\)};
    \node[above left] at (B) {\(e_2\)};
    \node[above right] at (W) {\(w=(1,1,2)\)};
  \end{scope}
  \node[font=\bfseries] at (0.40,2.95) {\(P\)};
  \node at (0.40,-0.95) {four lattice points};

  \begin{scope}[xshift=5.8cm,
    x={(0.86625cm,-0.21cm)},
    y={(-0.49875cm,0.1575cm)},
    z={(0.0525cm,0.5775cm)}]
    \coordinate (OO) at (0,0,0);
    \coordinate (AA) at (2,0,0);
    \coordinate (BB) at (0,2,0);
    \coordinate (WW) at (2,2,4);
    \coordinate (H) at (1,1,1);
    \coordinate (Qzero) at (1,1,0);
    \coordinate (Qtwo) at (1,1,2);
    \fill[PredictionTeal!7] (OO)--(AA)--(WW)--cycle;
    \fill[PredictionTeal!13] (OO)--(BB)--(WW)--cycle;
    \draw[PredictionTeal,thick,dashed] (AA)--(BB);
    \draw[PredictionTeal,very thick] (OO)--(AA)--(WW)--(BB)--cycle;
    \draw[PredictionTeal,thick] (OO)--(WW);
    \draw[black!50,densely dotted] (Qzero)--(Qtwo);
    % All six edge midpoints, including (2,1,2) and (1,2,2).
    \foreach \point in {(1,0,0),(0,1,0),(1,1,0),
                          (1,1,2),(2,1,2),(1,2,2)}
      \fill[black!75] \point circle (1.8pt);
    \foreach \point in {OO,AA,BB,WW}
      \fill[PredictionTeal] (\point) circle (2pt);
    \fill[white] (H) circle (3.6pt);
    \filldraw[fill=ExampleGold!28,draw=ExampleGold,thick]
      (H) circle (2.7pt);
    \node[below left] at (OO) {\(0\)};
    \node[below right] at (AA) {\(2e_1\)};
    \node[above left] at (BB) {\(2e_2\)};
    \node[above right] at (WW) {\(2w\)};
  \end{scope}
  \draw[ExampleGold,thin] (6.28,0.54)--(7.32,0.85);
  \node[right,text=ExampleGold,fill=white,inner sep=1.5pt]
    at (7.32,0.85) {\(h=(1,1,1)\)};
  \node[font=\bfseries] at (6.20,2.95) {\(2P\)};
  \node at (6.20,-0.95) {eleven lattice points};
\end{tikzpicture}
\caption{The height-two Reeve tetrahedron and its dilation, drawn at
matching display sizes. The only lattice points of \(P\) are its four
vertices. The eleven lattice points of \(2P\) consist of four vertices,
six edge midpoints, and the highlighted interior point \(h\), which
supplies the missing character \(\chi^h=stu\). Dashed edges are hidden.}
\label{fig:classroom-oct19-reeve}
\end{figure}

Now replace \(P\) by \(2P\), keeping the same normal fan. If \(V\)
denotes the set of four vertices of \(P\), then
\[
2P\cap M
=\{v+v':v,v'\in V\}\cup\{h\}.
\]
There are eleven points: four vertices, six edge midpoints, and the
interior point \(h\). They give the map to \(\PP^{10}\) with
coordinates
\[
[1:s:t:st:stu:stu^2:s^2:t^2:s^2tu^2:st^2u^2:s^2t^2u^4].
\]
The new coordinates recover \(u=(stu)/(st)\) on the torus. If the
coordinates corresponding to \(st,stu,stu^2\) are \(X_3,X_4,X_5\),
then \(X_4^2=X_3X_5\); this is one relation, rather than a complete
set of defining equations.

To check the boundary as well, let \(v\in V\), and let
\(r_1,r_2,r_3\) be its differences to the other three vertices. Their
determinant has absolute value two. The half-open fundamental
parallelepiped contains precisely \(0\) and
\[
h_v=h-2v=\tfrac12(r_1+r_2+r_3).
\]
Subtracting integer parts of the cone coefficients gives
\[
\Cone{P-v}\cap M
=\NN r_1+\NN r_2+\NN r_3+\NN h_v.
\]
All four generators occur in \((2P\cap M)-2v\): use the points
\(v+v_i\) and \(h\). The vertex criterion therefore proves that
\(2D_P\) is very ample. Doubling supplies the missing local
generators; it changes neither the fan nor the singularities.
\end{example}
% END CLASSROOM ADDITION: notes/20261019.tex

\subsection{What changes the normal fan?}

The normal fan records directions, not lengths or absolute position.
Consequently:

\begin{itemize}
\item translating \(P\) by a vector does not change \(\Sigma_P\);

\item multiplying \(P\) by a positive scalar does not change
\(\Sigma_P\);

\item changing an edge length without changing the directions of the
facets often preserves the same normal fan.
\end{itemize}

For example, every axis-parallel rectangle has the normal fan of
\[
\PP^1\times\PP^1,
\]
and every positive dilation of the standard triangle has the normal
fan of \(\PP^2\).

The normal fan changes when a facet disappears, a new facet becomes
active, or the polytope loses dimension. These are exactly the
phenomena visible in \cref{fig:blowup-convexity-positivity}.

\subsection{The normal fan of a divisor polytope}

Let \(D\) be a basepoint-free torus-invariant Cartier divisor on a
complete toric variety \(X_\Sigma\). Its Cartier data
\[
\{m_\sigma\}_{\sigma\in\Sigma(n)}
\]
lie in \(P_D\), and the distinct values among them are precisely the
vertices of \(P_D\).

\begin{theorem}[The fan remembered by a divisor]
\label{thm:normal-fan-divisor-polytope}
Let \(D\) be basepoint free. For every vertex \(v\) of \(P_D\), the
corresponding maximal cone of the normal fan is
\[
\sigma_v
=
\bigcup_{\substack{\sigma\in\Sigma(n)\\m_\sigma=v}}
\sigma.
\]
Consequently,
\[
\boxed{
\Sigma\text{ refines }\Sigma_{P_D}.
}
\]
In words: the normal fan of \(P_D\) is obtained by gluing together all
maximal cones that carry the same Cartier datum.
\end{theorem}

\noindent
\textit{Reference.}
This is Cox--Little--Schenck, Proposition~6.2.5.

\begin{proof}[Conceptual proof sketch]
For \(u\in\sigma\), convexity gives
\[
\psi_D(u)
=
\langle m_\sigma,u\rangle
=
\min_{m\in P_D}\langle m,u\rangle.
\]
Thus every direction in \(\sigma\) sees the Cartier datum
\(m_\sigma\) as a minimizing vertex of \(P_D\). If several adjacent
cones have the same datum, their union belongs to one normal cone.
If the data differ, the minimizing vertex changes, so one has crossed
a wall of the normal fan.
\end{proof}

Combining this theorem with strict convexity gives
\[
\boxed{
\begin{aligned}
D\text{ basepoint free}
&\Longrightarrow
\Sigma\text{ refines }\Sigma_{P_D},
\\[1mm]
D\text{ ample}
&\Longleftrightarrow
\Sigma=\Sigma_{P_D}.
\end{aligned}
}
\]
Strict convexity means that neighboring maximal cones never share the
same Cartier datum, so no wall is glued away.

\begin{figure}[ht]
\centering
\begin{tikzpicture}[scale=0.83]
  % 2H normal fan
  \begin{scope}[xshift=-5.2cm]
    \fill[DictionaryBlue!8] (0,0)--(2.2,0)--(0,2.2)--cycle;
    \fill[SurfaceGreen!8] (0,0)--(0,2.2)--(-1.6,-1.6)--cycle;
    \fill[ExampleGold!9] (0,0)--(-1.6,-1.6)--(2.2,0)--cycle;
    \draw[very thick,-{Latex[length=2mm]}] (0,0)--(2.55,0);
    \draw[very thick,-{Latex[length=2mm]}] (0,0)--(0,2.55);
    \draw[very thick,-{Latex[length=2mm]}] (0,0)--(-1.9,-1.9);
    \draw[WarningRed,dashed,very thick,-{Latex[length=2mm]}]
      (0,0)--(1.75,1.75);
    \node[WarningRed,above right] at (1.5,1.5) {forgotten ray};
    \node[font=\bfseries] at (0.25,3.05) {\(2H\)};
    \node[align=center,font=\small] at (0.25,-2.45)
      {normal fan of a triangle\\the exceptional wall is glued away};
  \end{scope}

  % 2H-E normal fan
  \begin{scope}
    \fill[DictionaryBlue!8]
      (0,0)--(2.2,0)--(1.5,1.5)--cycle;
    \fill[PredictionTeal!9]
      (0,0)--(1.5,1.5)--(0,2.2)--cycle;
    \fill[SurfaceGreen!8]
      (0,0)--(0,2.2)--(-1.6,-1.6)--cycle;
    \fill[ExampleGold!9]
      (0,0)--(-1.6,-1.6)--(2.2,0)--cycle;
    \draw[very thick,-{Latex[length=2mm]}] (0,0)--(2.55,0);
    \draw[very thick,-{Latex[length=2mm]}] (0,0)--(1.9,1.9);
    \draw[very thick,-{Latex[length=2mm]}] (0,0)--(0,2.55);
    \draw[very thick,-{Latex[length=2mm]}] (0,0)--(-1.9,-1.9);
    \node[font=\bfseries] at (0.25,3.05) {\(2H-E\)};
    \node[align=center,font=\small] at (0.25,-2.45)
      {the normal fan equals\\the blowup fan};
  \end{scope}

  % 2H-2E generalized normal fan
  \begin{scope}[xshift=5.2cm]
    \fill[SurfaceGreen!10]
      (-2.2,-2.2)--(2.2,2.2)--(2.2,-2.2)--cycle;
    \fill[ExampleGold!10]
      (-2.2,-2.2)--(2.2,2.2)--(-2.2,2.2)--cycle;
    \draw[very thick,PredictionTeal]
      (-2.35,-2.35)--(2.35,2.35);
    \draw[->,black!55] (-2.6,0)--(2.6,0);
    \draw[->,black!55] (0,-2.6)--(0,2.6);
    \node[PredictionTeal,rotate=45,fill=white,inner sep=1pt]
      at (1.25,1.4) {common lineality};
    \node[font=\bfseries] at (0.25,3.05) {\(2H-2E\)};
    \node[align=center,font=\small] at (0.25,-3.0)
      {a segment gives two half-planes;\\after quotienting, the fan is that of \(\PP^1\)};
  \end{scope}
\end{tikzpicture}
\caption{The normal fans associated to the three polytopes in
\cref{fig:blowup-convexity-positivity}. A redundant inequality glues
away the exceptional wall; an active facet remembers it; a
one-dimensional polytope produces a degenerate normal fan with
lineality.}
\label{fig:three-normal-fans-blowup}
\end{figure}

\subsection{The morphism encoded by the normal fan}

Because \(\Sigma\) refines \(\Sigma_{P_D}\), the identity on the
lattice induces a proper toric morphism
\[
\pi_D:X_\Sigma\longrightarrow X_{\Sigma_{P_D}}.
\]
When \(P_D\) is lower dimensional, one first quotients by the common
lineality space of its normal cones. The resulting toric variety has
dimension \(\dim P_D\).

\begin{theorem}[A basepoint-free divisor comes from an ample one]
\label{thm:basepoint-free-pullback-ample}
Let \(D\) be basepoint free on the complete toric variety
\(X_\Sigma\). Then \(D\) is linearly equivalent to the pullback of an
ample divisor on the projective toric variety associated to \(P_D\):
\[
D\sim\pi_D^*A_{P_D}.
\]
\end{theorem}

\noindent
\textit{Reference.}
This is Cox--Little--Schenck, Theorem~6.2.8.

\begin{proof}[Conceptual proof sketch]
The support function of \(D\) is
\[
\psi_D(u)
=
\min_{m\in P_D}\langle m,u\rangle.
\]
It is linear precisely on the cones of the normal fan of \(P_D\), so
it descends from \(\Sigma\) to \(\Sigma_{P_D}\). On the normal fan it
is strictly convex, hence defines an ample divisor \(A_{P_D}\). Pulling
that support function back along the refinement recovers
\(\psi_D\), and therefore recovers \(D\) up to linear equivalence.
\end{proof}

The three blowup examples now have a geometric interpretation:
\[
\begin{array}{c|c|c}
D
&
\text{normal fan of }P_D
&
\text{behavior of the projective map}
\\
\hline
2H
&
\text{fan of }\PP^2
&
\text{contracts the exceptional curve}
\\
2H-E
&
\text{fan of }\Blp{\PP^2}{p}
&
\text{forgets no toric wall}
\\
2H-2E
&
\text{fan of }\PP^1\text{ after quotient}
&
\text{records the pencil of lines through }p.
\end{array}
\]

\begin{surfacelaboratory}
Use \cref{fig:three-normal-fans-blowup} to explain the following maps.

\begin{enumerate}[label=(\alph*),wide=0pt]
\item Why does the map defined by \(2H\) factor through the blowdown
\[
\Blp{\PP^2}{p}\longrightarrow\PP^2?
\]

\item Why does \(2H-E\) retain the exceptional wall?

\item Why does \(2H-2E\) have one-dimensional image, even though its
complete linear system has three monomial sections?
\end{enumerate}
\end{surfacelaboratory}

A full-dimensional polytope and its normal fan determine each other
only up to changes of support numbers: many different polytopes can
have the same normal fan. On a toric surface, varying these support
numbers changes edge lengths while preserving the directions of the
boundary divisors.

\begin{guidingquestion}
Given the rays
\[
u_1,\dots,u_r
\]
of a complete toric surface, which polygons have these vectors as
primitive inward facet normals? How do their edge lengths and lattice
points encode sections and intersection geometry?
\end{guidingquestion}
% END INLINED: chapter05-section04-complete.tex

% BEGIN INLINED: chapter05-section05-complete.tex
\section{Toric surfaces through their polygons}
\label{sec:toric-surfaces-polygons}
\index{toric surface!projectivity}
\index{smooth polytope}
\index{intersection number!edge length}
\index{self-intersection!polytope area}
\index{corner cut!toric blowup}

In dimension two, the fan--polytope dictionary is especially complete.
A polygon can be drawn, measured, and modified by hand, while still
remembering the projective geometry of its toric surface.

The key dimension-two fact is the following.

\begin{theorem}[Every complete toric surface is projective]
\label{thm:complete-toric-surface-projective}
Every complete toric surface is projective. Equivalently, every
complete fan in \(N_\RR\simeq\RR^2\) is the normal fan of a lattice
polygon.
\end{theorem}

\noindent
\textit{Reference.}
This is Cox--Little--Schenck, Proposition~6.3.25.

\begin{proof}[Conceptual proof sketch]
Let
\[
u_1,\ldots,u_s
\]
be primitive generators of the rays, listed cyclically. Because the
fan is complete, the origin lies in the interior of their convex hull.
Hence there is a relation
\[
\lambda_1u_1+\cdots+\lambda_su_s=0
\]
with every \(\lambda_i>0\). The relation can be chosen rational and
then scaled so that all \(\lambda_i\) are positive integers.

Rotate each vector \(u_i\) by ninety degrees, choosing the direction
compatible with the cyclic ordering, and use the resulting vector of
lattice length \(\lambda_i\) as the \(i\)-th edge of a polygon. The
relation above says that the edge path closes. Its cyclic ordering
makes the polygon convex, and its primitive inward facet normals are
exactly the \(u_i\). Therefore its normal fan is the original fan.
\end{proof}

Thus every complete toric surface can be studied through many
polygons: changing the support numbers changes the divisor while
preserving the surface.

\subsection{A four-polygon catalogue}

The same familiar surfaces now appear as familiar shapes.

\begin{figure}[ht]
\centering
\begin{tikzpicture}[scale=0.72]
  % P2
  \begin{scope}[xshift=-6.1cm,yshift=2.45cm]
    \draw[->,black!50] (-0.25,0)--(2.75,0);
    \draw[->,black!50] (0,-0.25)--(0,2.75);
    \filldraw[fill=DictionaryBlue!13,draw=DictionaryBlue,very thick]
      (0,0)--(2.2,0)--(0,2.2)--cycle;
    \node[font=\bfseries] at (1.05,3.05) {\(\PP^2\)};
    \node[font=\small,align=center] at (1.05,-0.65)
      {triangle\\\(P_{dH}=d\Delta_2\)};
  \end{scope}

  % P1xP1
  \begin{scope}[xshift=0cm,yshift=2.45cm]
    \draw[->,black!50] (-0.25,0)--(2.95,0);
    \draw[->,black!50] (0,-0.25)--(0,2.75);
    \filldraw[fill=SurfaceGreen!13,draw=SurfaceGreen,very thick]
      (0,0)--(2.45,0)--(2.45,1.85)--(0,1.85)--cycle;
    \node[font=\bfseries] at (1.2,3.05)
      {\(\PP^1\times\PP^1\)};
    \node[font=\small,align=center] at (1.2,-0.65)
      {rectangle\\\(P_{aF_1+bF_2}\)};
  \end{scope}

  % Hirzebruch
  \begin{scope}[xshift=-6.1cm,yshift=-3.45cm]
    \draw[->,black!50] (-0.25,0)--(3.3,0);
    \draw[->,black!50] (0,-0.25)--(0,2.75);
    \filldraw[fill=ExampleGold!15,draw=ExampleGold,very thick]
      (0,0)--(1.75,0)--(2.8,1.9)--(0,1.9)--cycle;
    \node[font=\bfseries] at (1.35,3.05) {\(\mathbb F_r\)};
    \node[font=\small,align=center] at (1.35,-0.65)
      {trapezoid\\the slope records \(r\)};
  \end{scope}

  % blowup
  \begin{scope}[xshift=0cm,yshift=-3.45cm]
    \draw[->,black!50] (-0.25,0)--(2.95,0);
    \draw[->,black!50] (0,-0.25)--(0,2.75);
    \filldraw[fill=PredictionTeal!14,draw=PredictionTeal,very thick]
      (0.72,0)--(2.35,0)--(0,2.35)--(0,0.72)--cycle;
    \node[font=\bfseries] at (1.15,3.05)
      {\(\Blp{\PP^2}{p}\)};
    \node[font=\small,align=center] at (1.15,-0.65)
      {truncated triangle\\\(P_{dH-\ell E}\)};
  \end{scope}
\end{tikzpicture}
\caption{Four standard projective toric surfaces and characteristic
shapes of their divisor polygons. The normal directions determine the
surface; the edge positions and lengths record the chosen divisor.}
\label{fig:four-toric-surface-polygons}
\end{figure}

For the Hirzebruch surface, use the ray generators
\[
u_1=e_1,
\qquad
u_2=e_2,
\qquad
u_3=-e_1+re_2,
\qquad
u_4=-e_2.
\]
The divisor
\[
D=aD_3+bD_4
\]
has polygon
\[
P_D
=
\left\{
(x,y):
0\leq y\leq b,
\quad
0\leq x\leq a+ry
\right\},
\]
with vertices
\[
(0,0),
\quad
(a,0),
\quad
(a+rb,b),
\quad
(0,b).
\]
The parameter \(r\) is visible as the slope of the nonvertical side.

\begin{pauseandthink}
Starting only from this trapezoid, recover its four primitive inward
facet normals. In what order do they appear in the fan of
\(\mathbb F_r\)?
\end{pauseandthink}

\subsection{Smoothness is visible at the corners}

Let \(P\) be a lattice polygon. At each vertex, two facets meet, with
primitive inward normals \(u,v\in N\).

\begin{proposition}[Smooth polygon criterion]
\label{prop:smooth-polygon-criterion}
The projective toric surface \(X_{\Sigma_P}\) is smooth if and only if,
at every vertex of \(P\), the two primitive inward facet normals form
a basis of \(N\):
\[
|\det(u,v)|=1.
\]
Equivalently, the primitive directions of the two edges leaving each
vertex form a basis of \(M\).
\end{proposition}

\begin{proof}
A vertex of \(P\) corresponds to a maximal cone of \(\Sigma_P\), and
that cone is generated by the primitive inward normals of the two
incident facets. Apply the smoothness criterion for affine toric
charts.
\end{proof}

A lattice polygon satisfying this condition is often called a
\emph{smooth} or \emph{Delzant} polygon.

% BEGIN CLASSROOM ADDITION: notes/20261019.tex (primitive directions in Diagram 8).
\begin{observation}[Use primitive edge directions]
At \((2,0)\) in the polygon of
Example~\ref{ex:classroom-trapezoid-charts}, the primitive outgoing
directions are \((-1,0)\) and \((-1,1)\), with determinant \(-1\).
The entire bottom-edge vector is \((-2,0)\); using it would
incorrectly multiply the determinant by two. The smoothness test
records directions in the lattice, independently of edge lengths.
\end{observation}
% END CLASSROOM ADDITION: notes/20261019.tex

\begin{example}[A singular triangle]
The triangle
\[
P=\Conv{(0,0),(n,0),(0,1)}
\]
has inward normals
\[
e_1,
\qquad
e_2,
\qquad
-e_1-ne_2.
\]
At one vertex, the determinant of the two incident normals has
absolute value \(n\). Its normal fan therefore gives
\(\PP(1,1,n)\), which is singular for \(n>1\).
\end{example}

\subsection{Edge lengths are intersection numbers}

\begin{agcheatsheet}
Let \(X\) be a smooth complete surface, let \(D\) be a Cartier
divisor, and let \(C\subseteq X\) be an irreducible complete curve.
The intersection number
\[
D\cdot C
\]
is the degree of the restricted line bundle
\[
\mathcal O_X(D)|_C.
\]

If \(D\) gives a projective embedding, then \(D\cdot C\) is the degree
of the image of \(C\) with respect to a hyperplane. The
self-intersection \(D^2\) is the degree of the embedded surface.
\end{agcheatsheet}

% BEGIN CLASSROOM ADDITION: notes/20261019.tex (the geometric restriction picture).
\begin{observation}[Count zeros after restricting the bundle]
\label{obs:classroom-oct19-restrict-section}
Suppose \(C\) is smooth and a nonzero regular section
\(\widetilde s\) of \(L=\mathcal O_X(D)\) has zero divisor \(A\).
If \(C\) is not a component of \(A\), then
\[
D\cdot C=\deg(L|_C)
=\sum_{p\in C}\ord_p(\widetilde s|_C).
\]
Geometrically, intersect the section graph with the zero section,
then restrict to \(C\). Transverse intersections contribute one;
all intersections are counted with multiplicity.

If \(C\) is a component of \(A\), the restricted section is
identically zero, although the restricted line bundle still exists.
Choose another section when the linear system permits it; basepoint
freeness guarantees a section not vanishing identically on \(C\).
Any two nonzero rational sections differ by a rational function:
\[
\widetilde s'=g\widetilde s,\qquad
\operatorname{div}(\widetilde s')
=A+\operatorname{div}(g).
\]
This is the section interpretation of changing the divisor by a
principal divisor. The replacement is regular exactly when its
section divisor is effective. Without such a regular replacement,
one can use a rational section avoiding \(C\) at its generic point;
its restriction computes degree as zeros \emph{minus} poles. Thus
moving a divisor does not require pretending every line bundle is
basepoint free.
\end{observation}

\begin{figure}[htbp]
\centering
\begin{tikzpicture}[font=\small,>=Latex]
  \begin{scope}[xshift=-3.5cm]
    \filldraw[fill=DictionaryBlue!3,draw=black!35]
      (-2.3,0.3) rectangle (2.3,2.3);
    \draw[DictionaryBlue,thick] (-2.3,1)--(2.3,1)
      node[right] {\(0_C\)};
    \draw[WarningRed,very thick,domain=-2.1:2.1,samples=60]
      plot (\x,{1+0.25*(\x*\x-1)});
    \node[WarningRed] at (0,1.85) {\(\widetilde s|_C\)};
    \draw[thick] (-2.3,-0.4)--(2.3,-0.4) node[right] {\(C\)};
    \foreach \p in {-1,1} {
      \draw[dashed,black!45] (\p,1)--(\p,-0.4);
      \fill (\p,1) circle (1.8pt);
      \fill (\p,-0.4) circle (1.8pt);
    }
    \node[above] at (0,2.3) {\(L|_C\): isolated zeros};
  \end{scope}
  \begin{scope}[xshift=3.5cm]
    \filldraw[fill=DictionaryBlue!3,draw=black!35]
      (-2.3,0.3) rectangle (2.3,2.3);
    \draw[DictionaryBlue,line width=2pt] (-2.3,1)--(2.3,1);
    \draw[WarningRed,very thick,dashed] (-2.3,1)--(2.3,1);
    \node[WarningRed] at (0,1.85) {\(\widetilde s|_C\equiv0\)};
    \draw[thick] (-2.3,-0.4)--(2.3,-0.4) node[right] {\(C\)};
    \node[above] at (0,2.3) {a component of the zero divisor};
  \end{scope}
\end{tikzpicture}
\caption{Schematic slices of the restricted line bundle. On the left,
the section gives a finite zero divisor on the curve. On the right,
it coincides with the zero section, so a different rational or regular
representative is needed to compute degree by points.}
\label{fig:classroom-oct19-bundle-restriction}
\end{figure}
% END CLASSROOM ADDITION: notes/20261019.tex

For an edge \(F\) of a lattice polygon, define its \emph{lattice
length} by
\[
\length_M(F)
=
|F\cap M|-1.
\]
Equivalently, if its endpoints differ by \(\ell w\), where \(w\) is
primitive, then \(\length_M(F)=\ell\).

\begin{theorem}[The edge--intersection dictionary]
\label{thm:edge-length-intersection}
Let \(X_\Sigma\) be a smooth projective toric surface, and let \(D\)
be an ample torus-invariant Cartier divisor. If \(F_\rho\) is the edge
of \(P_D\) with primitive inward normal \(u_\rho\), then
\[
\boxed{
D\cdot D_\rho
=
\length_M(F_\rho).
}
\]
\end{theorem}

\noindent
\textit{Reference.}
This is the surface interpretation of Cox--Little--Schenck,
Proposition~6.3.8.

\begin{proof}[Conceptual proof sketch]
The invariant curve \(D_\rho\) is a copy of \(\PP^1\). The two
vertices of \(F_\rho\) are the Cartier data of the two maximal cones
adjacent to \(\rho\). Their difference lies in the rank-one lattice
\[
M\cap u_\rho^\perp
\]
and is \(\ell\) times its primitive generator, where
\(\ell=\length_M(F_\rho)\). Proposition~6.3.8 computes the degree of
\(\mathcal O(D)|_{D_\rho}\) from exactly this difference, giving
\(D\cdot D_\rho=\ell\).
\end{proof}

For
\[
D=dH-\ell E
\quad\text{on}\quad
\Blp{\PP^2}{p},
\]
the four edge lengths of
\[
P_D=
\{(a,b):a,b\geq0,\ \ell\leq a+b\leq d\}
\]
are
\[
d,
\qquad
d-\ell,
\qquad
d-\ell,
\qquad\ell.
\]
They recover the intersections
\[
D\cdot H=d,
\qquad
D\cdot(H-E)=d-\ell,
\qquad
D\cdot E=\ell.
\]
Thus positivity on every invariant curve can literally be read as
positivity of every edge length.

\begin{surfacelaboratory}
For the Hirzebruch trapezoid
\[
\Conv{(0,0),(a,0),(a+rb,b),(0,b)},
\]
compute the four lattice edge lengths. Match them with the four
intersection numbers \(D\cdot D_i\).

Which lengths vanish when \(a=0\) or \(b=0\), and which toric curves
should the corresponding projective map contract?
\end{surfacelaboratory}

% BEGIN CLASSROOM ADDITION: notes/20261019.tex (Hirzebruch classes and restriction).
\subsection{Supplementary example: the classroom trapezoid}
\label{subsec:classroom-oct19-hirzebruch-restriction}

Return to the polygon from Example~\ref{ex:classroom-trapezoid-charts}, in
the same lattice coordinates:
\[
P=\Conv\{(0,0),(2,0),(1,1),(0,1)\}.
\]
Its inward normal rays are \(e_1,e_2,-e_1-e_2,-e_2\). Their fan is
the fan of \(X=\Hirz{1}\cong\Blp{\PP^2}{p}\). Name the four
invariant curves
\[
F=D_{e_1},\qquad B=D_{e_2},\qquad
F'=D_{-e_1-e_2},\qquad E=D_{-e_2}.
\]
The notation \(D_u\) here means the divisor of the ray generated by
\(u\). The coefficients of the divisor defining \(P\) are
\[
D=2F'+E,\qquad
P_D=\{(r,q):r,q\geq0,\ q\leq1,\ r+q\leq2\}.
\]
These coordinates differ from the earlier Hirzebruch catalogue:
the affine lattice transformation \((r,q)\mapsto(r,1-q)\) sends
this polygon to \(\Conv\{(0,0),(1,0),(2,1),(0,1)\}\).
Its linear part reflects the normal fan. In particular, the names
of the upper and lower section curves must be matched through that
reflection.

\begin{example}[Recover the intersections from the class group]
\label{ex:classroom-oct19-hirzebruch-classes}
Write \(s=\character{(1,0)}\) and \(t=\character{(0,1)}\).
The character-divisor formula gives
\[
\operatorname{div}(s)=F-F',\qquad
\operatorname{div}(t)=B-F'-E.
\]
Thus \(F'\sim F\) and \(B\sim E+F\). The divisor-class sequence
from Chapter~4 gives
\[
\Cl(X)=\Pic(X)=\ZZ[F]\oplus\ZZ[E].
\]
Here \(F,F'\) are distinct fibers of the ruling, \(E\) is its
negative section, and \(B\) is a disjoint section. Adjacent invariant
curves meet transversely once, while nonadjacent ones are disjoint.
It follows that
\[
F^2=F\cdot F'=0,\qquad E\cdot F=1,\qquad
0=E\cdot B=E\cdot(E+F),
\]
so \(E^2=-1\) and \(B^2=(E+F)^2=1\). The full intersection table is
\[
\begin{array}{c|rrrr}
\cdot&F&B&F'&E\\ \hline
F &0&1&0&1\\
B &1&1&1&0\\
F'&0&1&0&1\\
E &1&0&1&-1
\end{array}
\]
Taking \(H=B\), the pullback of a line on \(\PP^2\), recovers
\[
D\sim2F+E=2H-E.
\]
Thus the class-group calculation gives exactly the polygon lengths:
\[
\begin{array}{c|c|c}
\text{curve}&\text{dual edge of }P_D&D\cdot\text{curve}\\ \hline
F &r=0&1\\
B &q=0&2\\
F'&r+q=2&1\\
E &q=1&1
\end{array}
\]
For instance, \(F\cdot(2F+E)=1\), whereas
\(B\cdot(2F+E)=2\). The four numbers come from the same divisor
class, even though its chosen invariant representative visibly
contains only \(F'\) and \(E\).
\end{example}

\begin{example}[Restrict using the star fan]
\label{ex:classroom-oct19-star-restriction}
The invariant curve \(C=V(\rho)\) is the \emph{closure} of the orbit
\(O(\rho)\cong\Gm\), and here \(C\cong\PP^1\). By
Proposition~\ref{prop:fan-of-orbit-closure}, its fan is formed by retaining only
cones containing \(\rho\), then taking their images in
\[
N_\rho=N/\ZZ u_\rho.
\]
One must restrict to this star before taking the quotient; projecting
every ray can include a divisor that does not meet \(C\).

First let \(C=F\). Its coefficient in \(D=2F'+E\) is zero, so
the canonical section does not vanish identically on it. The
neighboring rays are \(e_2\) and \(-e_2\), and the quotient map
\(q(a,b)=b\) sends them to \(+1\) and \(-1\). Their coefficients
are \(0\) and \(1\). Hence
\[
D|_F=[E\cap F],\qquad
\mathcal O_X(D)|_F\cong\mathcal O_{\PP^1}(1).
\]
Although \(-e_1-e_2\) also maps to \(-1\), its divisor \(F'\) is
disjoint from \(F\) and contributes nothing.

% Classroom addition: notes/20261019.tex, Diagram 11 (restriction to a fiber).
\begin{figure}[htbp]
\centering
\begin{tikzpicture}[font=\small,>={Latex[length=2.2mm]},line join=round]
  % Retain only the two maximal cones containing rho_F.
  \fill[SurfaceGreen!9] (0,0) rectangle (1.75,1.75);
  \fill[PausePurple!9] (0,0) rectangle (1.75,-1.75);
  \draw[very thick,DictionaryBlue,->] (0,0)--(1.95,0)
    node[right] {\(F:e_1\)};
  \draw[very thick,SurfaceGreen,->] (0,0)--(0,1.95)
    node[above] {\(B:e_2\)};
  \draw[very thick,PausePurple,->] (0,0)--(0,-1.95)
    node[below] {\(E:-e_2\)};
  \draw[thick,WarningRed!55!black,densely dashed,->]
    (0,0)--(-1.35,-1.35)
    node[below left] {\(F'\)};
  \fill (0,0) circle (1.5pt);
  \node[DictionaryBlue,below] at (1.05,-0.10) {\(a_F=0\)};
  \node[SurfaceGreen,left] at (-0.12,1.15) {\(a_B=0\)};
  \node[PausePurple,left] at (-0.12,-1.50) {\(a_E=1\)};
  \node[WarningRed!65!black,left] at (-0.55,-0.48) {\(a_{F'}=2\)};
  \node at (0,2.75) {\(\Sigma_P\)};
  \node[align=center] at (0,-2.78)
    {\(D=2F'+E\)\\shaded: cones containing \(\rho_F\)};

  \draw[thick,->] (3.25,0.18)--(5.10,0.18)
    node[midway,above] {\(q(a,b)=b\)}
    node[midway,below] {\(e_1\mapsto0\)};

  \begin{scope}[xshift=6.55cm]
    \draw[very thick,SurfaceGreen,->] (0,0)--(0,1.90)
      node[above] {\(+1\)};
    \draw[very thick,PausePurple,->] (0,0)--(0,-1.90)
      node[below] {\(-1\)};
    \fill (0,0) circle (1.5pt);
    \node[left] at (-0.12,0) {\(0\)};
    \node[SurfaceGreen,right] at (0.15,1.12) {\(a_B=0\)};
    \node[PausePurple,right] at (0.15,-1.12) {\(a_E=1\)};
    \node at (0,2.75) {\(\operatorname{Star}(\rho_F)\)};
    \node[align=center] at (0,-2.78)
      {\(N/\ZZ e_1\cong\ZZ\)\\\(D|_F=[E\cap F]\)};
  \end{scope}
\end{tikzpicture}
\caption{Restrict first to the cones containing \(\rho_F\), then quotient
by \(\ZZ e_1\). The resulting fan describes the orbit closure
\(F=V(\rho_F)\cong\PP^1\), whose dense orbit is \(O(\rho_F)\cong\Gm\).
Since \(a_F=0\), the adjacent coefficients \(0,1\) descend directly,
giving a divisor of degree one. The dashed ray of the disjoint fiber
\(F'\) is outside this star and contributes nothing.}
\label{fig:classroom-oct19-star-restriction}
\end{figure}

Next take \(C=F'\). The canonical section now vanishes identically
on \(C\), since its coefficient in \(D\) is two. The characters
\(st\) and \(s^2\) give two regular replacements:
\[
\begin{aligned}
A_1&=D+\operatorname{div}(st)=F+B,\\
A_2&=D+\operatorname{div}(s^2)=2F+E.
\end{aligned}
\]
Both are effective and have coefficient zero on \(F'\). This is the
board's instruction to add a linear function. In the chapter's
support-function convention the precise formula is
\[
\psi_{D+\operatorname{div}(\character{m})}(u)
=\psi_D(u)-\langle m,u\rangle.
\]
Equivalently, the opposite function \(\eta_D=-\psi_D\) from
Observation~\ref{obs:classroom-support-sign} gains \(x+y\) or \(2x\).
Both have value \(-2\) at \((-1,-1)\), cancelling the coefficient
two there. The new linear pieces on the star annihilate that ray
and descend to the quotient.

For \(F'\), use \(q'(a,b)=b-a\). The neighboring rays of \(B\)
and \(E\) map to \(+1\) and \(-1\); the other fiber \(F\) is
disjoint. With \(p_B=B\cap F'\), \(p_E=E\cap F'\), we obtain
\[
A_1|_{F'}=p_B,\qquad A_2|_{F'}=p_E.
\]
The coefficient pairs \((1,0)\) and \((0,1)\) both represent
\(\mathcal O_{\PP^1}(1)\). They differ by the principal divisor
of the quotient character \(z=(t/s)|_{F'}\):
\[
\operatorname{div}_{F'}(z)=p_B-p_E.
\]
This explains why different choices of moving section give the same
intersection number.

On the polygon side, the edge dual to \(F'\) is
\(G=[(2,0),(1,1)]\). Translating by the two character exponents puts
it into the character lattice of the curve:
\[
M\cap(-1,-1)^\perp=\ZZ(-1,1).
\]
In that lattice,
\[
\begin{aligned}
G-(1,1)&=[(1,-1),(0,0)]\cong[-1,0],\\
G-(2,0)&=[(0,0),(-1,1)]\cong[0,1].
\end{aligned}
\]
The two intervals differ by a lattice translation and both have
length one, recovering \(D\cdot F'=1\).
\end{example}

\begin{observation}[Why a negative self-intersection needs a pole]
\label{obs:classroom-oct19-negative-section}
Apply the same procedure to \(\mathcal O_X(E)|_E\). Since
\[
E+\operatorname{div}(t)=B-F',
\]
a rational replacement has coefficient zero on \(E\). Restricting
gives \(-[E\cap F']\), because \(B\) is disjoint from \(E\).
There is one pole and no zero, so the degree is \(-1\). This is
\(E^2=-1\); no nonzero regular section of the restricted line bundle
could have that degree.
\end{observation}
% END CLASSROOM ADDITION: notes/20261019.tex

\subsection{Area is self-intersection}

Normalize area so that a fundamental lattice parallelogram in
\(M\) has area \(1\).

\begin{theorem}[Area and self-intersection]
\label{thm:polygon-area-self-intersection}
Let \(D\) be a nef torus-invariant Cartier divisor on a smooth complete
toric surface. Then
\[
\boxed{
D^2=2\operatorname{Area}(P_D).
}
\]
Moreover,
\[
-D\cdot K_X
=
|\partial P_D\cap M|.
\]
\end{theorem}

\noindent
\textit{Reference.}
This is Cox--Little--Schenck, Proposition~10.5.6; see also
Theorem~13.4.1.

\begin{proof}[Conceptual proof sketch]
The number of sections of \(kD\) is
\[
|kP_D\cap M|.
\]
As \(k\) grows, the leading term of this lattice-point count is
\[
\operatorname{Area}(P_D)k^2.
\]
On the geometric side, the leading term of the Hilbert polynomial is
\[
\frac{D^2}{2}k^2.
\]
Comparing the leading coefficients gives the area formula. Comparing
the linear coefficients gives the boundary formula.
\end{proof}

% BEGIN CLASSROOM ADDITION: notes/20261019.tex (degree of the classroom embedding).
For the classroom trapezoid in
\cref{subsec:classroom-oct19-hirzebruch-restriction}, the area is
\(3/2\). Thus
\[
D^2=2\operatorname{Area}(P_D)=3,
\]
also obtained from \((2F+E)^2=4(F\cdot E)+E^2=3\).
The embedding by its five sections therefore realizes
\(\Hirz{1}\) as a degree-three surface in \(\PP^4\).
% END CLASSROOM ADDITION: notes/20261019.tex

When \(P_D\) is two-dimensional, combining these identities with
Pick's theorem,
\[
\operatorname{Area}(P)
=
|\operatorname{Int}(P)\cap M|
+
\frac12|\partial P\cap M|-1,
\]
produces the surface Riemann--Roch count
\[
|P_D\cap M|
=
\frac{D^2-D\cdot K_X}{2}+1.
\]
In toric geometry, the lattice-point formula and the intersection
formula are two views of the same polynomial.

\begin{observation}
For an ample divisor, the interior lattice points also have a direct
divisor meaning. Since
\[
K_X=-\sum_\rho D_\rho,
\]
we have
\[
P_{D+K_X}\cap M
=
\operatorname{Int}(P_D)\cap M.
\]
Thus interior lattice points index the monomial sections of
\(\mathcal O(D+K_X)\).
\end{observation}

\subsection{Blowing up cuts a corner}

Suppose \(P\) is smooth at a vertex \(v\), and the two incident facets
have primitive inward normals \(u\) and \(v'\). In suitable integral
coordinates near that corner, the polygon is locally
\[
x\geq0,
\qquad
y\geq0.
\]
Cutting off a corner of lattice size \(\ell\) adds the inequality
\[
x+y\geq\ell.
\]
The new facet has primitive inward normal
\[
u+v'.
\]
This is precisely the star subdivision that performs the toric blowup
of the corresponding fixed point.

\begin{figure}[ht]
\centering
\begin{tikzpicture}[scale=0.9]
  % before
  \begin{scope}[xshift=-4.0cm]
    \draw[->,black!55] (-0.25,0)--(3.0,0) node[below] {\(x\)};
    \draw[->,black!55] (0,-0.25)--(0,3.0) node[left] {\(y\)};
    \fill[DictionaryBlue!11] (0,0)--(2.55,0)--(2.55,2.55)--(0,2.55)--cycle;
    \draw[DictionaryBlue,very thick] (0,2.55)--(0,0)--(2.55,0);
    \fill (0,0) circle (2.4pt);
    \node[DictionaryBlue,left] at (0,1.35) {normal \(u\)};
    \node[DictionaryBlue,below] at (1.35,0) {normal \(v'\)};
    \node[font=\bfseries] at (1.1,3.45) {before the blowup};
    \node[font=\small] at (1.1,-0.7) {one smooth vertex};
  \end{scope}

  \draw[very thick,-{Latex[length=2.8mm]}]
    (-0.8,0.75)--(0.8,0.75);
  \node[above] at (0,0.75) {cut the corner};

  % after
  \begin{scope}[xshift=4.0cm]
    \draw[->,black!55] (-0.25,0)--(3.0,0) node[below] {\(x\)};
    \draw[->,black!55] (0,-0.25)--(0,3.0) node[left] {\(y\)};
    \fill[PredictionTeal!12]
      (0,0.75)--(0.75,0)--(2.55,0)--(2.55,2.55)--(0,2.55)--cycle;
    \draw[DictionaryBlue,very thick] (0,2.55)--(0,0.75);
    \draw[PredictionTeal,very thick]
      (0,0.75)--(0.75,0)
      node[midway,below,sloped,text=PredictionTeal,yshift=-4pt]
      {length \(\ell\)};
    \draw[DictionaryBlue,very thick] (0.75,0)--(2.55,0);
    \draw[PredictionTeal,->,thick]
      (0.38,0.38)--(1.12,1.12)
      node[above right,text=PredictionTeal,fill=white,inner sep=1pt]
      {normal \(u+v'\)};
    \node[font=\bfseries] at (1.1,3.45) {after the blowup};
    \node[font=\small] at (1.1,-0.7)
      {the new edge corresponds to \(E\)};
  \end{scope}
\end{tikzpicture}
\caption{At a smooth polygon corner, a toric blowup cuts off a
triangle. The new edge has inward normal equal to the sum of the two
old normals and corresponds to the exceptional divisor.}
\label{fig:polygon-corner-cut-blowup}
\end{figure}

Conversely, when an edge shrinks to length zero, the corresponding ray
disappears from the normal fan and the invariant curve is contracted.
This is the polygon version of a toric blowdown.

\begin{pauseandthink}
Begin with the triangle \(d\Delta_2\). Cut all three corners by lattice
sizes \(\ell_0,\ell_1,\ell_2\).

Write the six inequalities of the resulting polygon. What numerical
conditions must the \(\ell_i\) satisfy for all six edges to have
positive length?
\end{pauseandthink}

\subsection{A surface-recognition workflow}

For a lattice polygon \(P\), one can now read the associated surface
and divisor in a consistent order:

\begin{center}
\small
\renewcommand{\arraystretch}{1.3}
\begin{tabularx}{0.94\linewidth}{@{}
  >{\raggedright\arraybackslash}p{0.20\linewidth}
  >{\raggedright\arraybackslash}p{0.32\linewidth}
  >{\raggedright\arraybackslash}X@{}}
\toprule
\textbf{Look at}
&
\textbf{Combinatorial information}
&
\textbf{Geometric meaning}
\\
\midrule
facet normals
&
normal fan
&
the toric surface
\\
corners
&
determinants of adjacent normals
&
smooth or quotient-singular fixed points
\\
edge lengths
&
lattice lengths
&
intersection with invariant curves
\\
area
&
normalized two-dimensional volume
&
self-intersection of the divisor
\\
lattice points
&
monomial exponents
&
global sections and projective coordinates
\\
corner cuts
&
new facets and new normals
&
toric blowups and exceptional curves
\\
\bottomrule
\end{tabularx}
\end{center}

A polygon is therefore not merely a picture of a toric surface. It is
a simultaneous record of its fan, an ample divisor, its sections, its
projective map, and several intersection numbers.

The next chapter will replace the single projective coordinate system
coming from one divisor by a universal coordinate system with one
variable for every ray.

\begin{guidingquestion}
Can one build one polynomial ring
\[
\Bbbk[x_\rho:\rho\in\Sigma(1)]
\]
whose variables record all invariant boundary divisors at once, and
then recover \(X_\Sigma\) as a quotient of an open subset of affine
space?
\end{guidingquestion}
% END INLINED: chapter05-section05-complete.tex

% Chapter-end exercise: scaling, translation under principal divisors,
% and the Minkowski inclusion P_D+P_E\subseteq P_{D+E}.

\clearpage
\section*{What this chapter has accomplished}
\addcontentsline{toc}{section}{What this chapter has accomplished}
\markboth{What this chapter has accomplished}{What this chapter has accomplished}

A torus-invariant divisor
\[
D=\sum_{\rho\in\Sigma(1)}a_\rho D_\rho
\]
produces the polyhedron
\[
P_D
=
\left\{
 m\in M_\RR:
 \langle m,u_\rho\rangle\geq-a_\rho
 \text{ for every }\rho
\right\}.
\]
Its lattice points are precisely the exponents of its monomial global
sections:
\[
H^0(X_\Sigma,\mathcal O(D))
=
\bigoplus_{m\in P_D\cap M}\Bbbk\character{m}.
\]

Those sections define projective maps. For a finite set
\[
A=\{m_0,\ldots,m_r\}\subseteq M,
\]
the monomial map
\[
t\longmapsto
[\character{m_0}(t):\cdots:\character{m_r}(t)]
\]
produces a projective toric variety by taking the orbit closure.

The support function then detects how much projective geometry the
sections see:
\[
\boxed{
\begin{aligned}
\psi_D\text{ convex}
&\Longleftrightarrow
D\text{ is basepoint free},
\\
\psi_D\text{ strictly convex}
&\Longleftrightarrow
D\text{ is ample}.
\end{aligned}
}
\]
Equivalently, convexity makes the complete linear system define a
morphism, while strict convexity makes a sufficiently large multiple
give an embedding.

A polytope remembers these bends through its normal fan. For a
basepoint-free divisor,
\[
\Sigma\text{ refines }\Sigma_{P_D},
\]
and strict convexity is exactly the condition
\[
\Sigma=\Sigma_{P_D}.
\]
Thus an ample toric divisor is the same data as a lattice polytope
whose normal fan is the fan of the variety.

For surfaces, the dictionary becomes especially visible:
\[
\begin{array}{c@{\quad\longleftrightarrow\quad}c}
\text{edge lengths of }P_D
&
D\cdot D_\rho,
\\
2\operatorname{Area}(P_D)
&
D^2,
\\
P_D\cap M
&
\text{global sections},
\\
\text{cutting a corner}
&
\text{a toric blowup}.
\end{array}
\]
Every complete toric surface is projective. The challenge exercise
shows that this phenomenon stops in dimension three.

\begin{exercises}[Comprehension exercises]
\markboth{Comprehension exercises}{Comprehension exercises}

\item Let
\[
D=aF_1+bF_2
\]
on \(\PP^1\times\PP^1\), with \(a,b\geq0\).
\begin{enumerate}[label=(\alph*),wide=0pt]
\item Show that \(P_D\) is a rectangle with side lengths \(a\) and
\(b\).
\item Count its lattice points and recover
\[
h^0(\PP^1\times\PP^1,\mathcal O(a,b))=(a+1)(b+1).
\]
\item Write the monomial map associated to the complete linear system.
\item Determine when \(D\) is very ample and describe what the map does
when one of \(a,b\) is zero.
\end{enumerate}

\item Let \(D,E\) be torus-invariant divisors and let \(m_0\in M\).
Prove the following statements directly from the defining
inequalities:
\begin{enumerate}[label=(\alph*),wide=0pt]
\item
\[
P_{kD}=kP_D
\qquad(k\in\ZZ_{>0}).
\]
\item
\[
P_{D+\operatorname{div}(\character{m_0})}
=P_D-m_0.
\]
\item
\[
P_D+P_E\subseteq P_{D+E}.
\]
\end{enumerate}
Interpret each statement using sections and projective maps.

\item Let
\[
A=\{m_0,\ldots,m_r\}\subseteq M.
\]
\begin{enumerate}[label=(\alph*),wide=0pt]
\item Show that translating every element of \(A\) by the same lattice
vector does not change the monomial map to \(\PP^r\).
\item Show that the kernel of
\[
\phi_A:T\longrightarrow\PP^r
\]
is determined by the subgroup generated by the differences
\[
m_i-m_0.
\]
\item Give a criterion for \(\phi_A|_T\) to be injective.
\end{enumerate}

\item On \(X=\Blp{\PP^2}{p}\), consider
\[
D=dH-rE,
\qquad
0\leq r\leq d.
\]
\begin{enumerate}[label=(\alph*),wide=0pt]
\item Count the lattice points in
\[
P_D=\{(a,b):a,b\geq0,\ r\leq a+b\leq d\}.
\]
\item Recover the dimension of the space of degree-\(d\) plane curves
vanishing to order at least \(r\) at \(p\).
\item Determine when the complete linear system is basepoint free and
when \(D\) is ample.
\item Describe the maps arising in the boundary cases \(r=0\) and
\(r=d\).
\end{enumerate}

\item Let
\[
P=\Conv\{(0,0),(b,0),(b+ra,a),(0,a)\},
\qquad
a,b>0,
\]
where \(r\geq0\).
\begin{enumerate}[label=(\alph*),wide=0pt]
\item Compute the primitive inward normal to each edge.
\item Show that the normal fan is the fan of the Hirzebruch surface
\(\mathbb F_r\).
\item Determine the lattice lengths of the four edges.
\item Translate those lengths into intersection numbers of the
corresponding ample divisor with the invariant curves.
\end{enumerate}

\item Let \(P\subseteq M_\RR\) be a full-dimensional lattice polygon.
\begin{enumerate}[label=(\alph*),wide=0pt]
\item Show that \(P\) is smooth exactly when the primitive inward
normals to the two edges meeting at each vertex form a lattice basis of
\(N\).
\item Apply the test to the standard triangle, a rectangle, and the
trapezoid from the previous exercise.
\item Construct a lattice triangle whose normal fan gives
\(\PP(1,1,n)\), and identify its singular corner.
\end{enumerate}

\item For \(D=2H-E\) on \(\Blp{\PP^2}{p}\):
\begin{enumerate}[label=(\alph*),wide=0pt]
\item List the five lattice points of \(P_D\).
\item Write the associated monomial map on the dense torus.
\item Compute the lattice lengths of the four edges and recover the
intersection numbers of \(D\) with the four invariant curves.
\item Compute \(D^2\) from the area of \(P_D\).
\end{enumerate}

\item Let \(P\) be a smooth lattice polygon, and let \(P'\) be obtained
by cutting off a corner by a segment of lattice length \(r\).
\begin{enumerate}[label=(\alph*),wide=0pt]
\item Show that the new inward normal is the sum of the two old inward
normals.
\item Interpret the operation as a toric blowup.
\item Determine how the area changes.
\item Compare the change in area with
\[
(\pi^*D-rE)^2=D^2-r^2.
\]
\end{enumerate}

\item Use Pick's theorem to show that, for an ample divisor \(D\) on a
smooth projective toric surface,
\[
h^0(X,\mathcal O(D))
=
\operatorname{Area}(P_D)
+\frac12\,|\partial P_D\cap M|
+1.
\]
Then rewrite the right-hand side using \(D^2\), \(D\cdot(-K_X)\), and
\(\chi(\mathcal O_X)=1\).

\end{exercises}

\clearpage
\section*{Challenge exercise}
\addcontentsline{toc}{section}{Challenge exercise}
\markboth{Challenge exercise}{Challenge exercise}

\noindent
\textbf{Complete but not projective: Fulton's toric threefold.}

\medskip

For toric surfaces, completeness forces projectivity. We now explore
a smooth complete toric threefold for which no ample divisor exists.
The obstruction will be a cycle of strict-convexity inequalities that
cannot all hold simultaneously.

Let
\[
N=\ZZ^3
\]
with standard basis \(e_1,e_2,e_3\), and define
\[
\begin{aligned}
v_1&=-e_1,&
 v_2&=-e_2,&
 v_3&=-e_3,&
 v_4&=e_1+e_2+e_3,
\\
v_5&=v_3+v_4=e_1+e_2,&
 v_6&=v_1+v_4=e_2+e_3,&
 v_7&=v_2+v_4=e_1+e_3.
\end{aligned}
\]
Write
\[
\sigma_{ijk}=\Cone{v_i,v_j,v_k}.
\]
Consider the fan \(\Sigma\) whose maximal cones are
\[
\begin{aligned}
&\sigma_{123},
\\
&\sigma_{126},\ \sigma_{267},\ \sigma_{467},
\\
&\sigma_{135},\ \sigma_{156},\ \sigma_{456},
\\
&\sigma_{237},\ \sigma_{357},\ \sigma_{457}.
\end{aligned}
\]
Let
\[
X=X_\Sigma.
\]

\begin{figure}[ht]
\centering
\begin{tikzpicture}[scale=0.86]
  % Face 123
  \begin{scope}[xshift=-3.7cm,yshift=2.35cm]
    \coordinate (A) at (-1.15,0);
    \coordinate (B) at (1.15,0);
    \coordinate (C) at (0,1.85);
    \filldraw[fill=DictionaryBlue!10,draw=DictionaryBlue,very thick]
      (A)--(B)--(C)--cycle;
    \node[below left] at (A) {\(v_1\)};
    \node[below right] at (B) {\(v_2\)};
    \node[above] at (C) {\(v_3\)};
    \node[font=\bfseries] at (0,2.65) {face \(123\)};
  \end{scope}

  % Face 124
  \begin{scope}[xshift=3.7cm,yshift=2.35cm]
    \coordinate (A) at (-1.15,0);
    \coordinate (B) at (1.15,0);
    \coordinate (C) at (0,1.85);
    \coordinate (P) at (-0.56,0.94);
    \coordinate (Q) at (0.56,0.94);
    \fill[PredictionTeal!7] (A)--(B)--(P)--cycle;
    \fill[PredictionTeal!12] (B)--(P)--(Q)--cycle;
    \fill[PredictionTeal!18] (C)--(P)--(Q)--cycle;
    \draw[PredictionTeal,very thick]
      (A)--(B)--(Q)--(C)--(P)--cycle;
    \draw[PredictionTeal,very thick] (B)--(P)--(Q);
    \fill[PredictionTeal] (P) circle (2pt);
    \fill[PredictionTeal] (Q) circle (2pt);
    \node[below left] at (A) {\(v_1\)};
    \node[below right] at (B) {\(v_2\)};
    \node[above] at (C) {\(v_4\)};
    \node[left] at (P) {\(v_6\)};
    \node[right] at (Q) {\(v_7\)};
    \node[font=\bfseries] at (0,2.65) {face \(124\)};
  \end{scope}

  % Face 134
  \begin{scope}[xshift=-3.7cm,yshift=-2.0cm]
    \coordinate (A) at (-1.15,0);
    \coordinate (B) at (1.15,0);
    \coordinate (C) at (0,1.85);
    \coordinate (P) at (-0.56,0.94);
    \coordinate (Q) at (0.56,0.94);
    \fill[SurfaceGreen!7] (A)--(B)--(Q)--cycle;
    \fill[SurfaceGreen!12] (A)--(Q)--(P)--cycle;
    \fill[SurfaceGreen!18] (C)--(P)--(Q)--cycle;
    \draw[SurfaceGreen,very thick]
      (A)--(B)--(Q)--(C)--(P)--cycle;
    \draw[SurfaceGreen,very thick] (A)--(Q)--(P);
    \fill[SurfaceGreen] (P) circle (2pt);
    \fill[SurfaceGreen] (Q) circle (2pt);
    \node[below left] at (A) {\(v_1\)};
    \node[below right] at (B) {\(v_3\)};
    \node[above] at (C) {\(v_4\)};
    \node[left] at (P) {\(v_6\)};
    \node[right] at (Q) {\(v_5\)};
    \node[font=\bfseries] at (0,2.65) {face \(134\)};
  \end{scope}

  % Face 234
  \begin{scope}[xshift=3.7cm,yshift=-2.0cm]
    \coordinate (A) at (-1.15,0);
    \coordinate (B) at (1.15,0);
    \coordinate (C) at (0,1.85);
    \coordinate (P) at (-0.56,0.94);
    \coordinate (Q) at (0.56,0.94);
    \fill[ExampleGold!8] (A)--(B)--(P)--cycle;
    \fill[ExampleGold!13] (B)--(P)--(Q)--cycle;
    \fill[ExampleGold!19] (C)--(P)--(Q)--cycle;
    \draw[ExampleGold,very thick]
      (A)--(B)--(Q)--(C)--(P)--cycle;
    \draw[ExampleGold,very thick] (B)--(P)--(Q);
    \fill[ExampleGold] (P) circle (2pt);
    \fill[ExampleGold] (Q) circle (2pt);
    \node[below left] at (A) {\(v_2\)};
    \node[below right] at (B) {\(v_3\)};
    \node[above] at (C) {\(v_4\)};
    \node[left] at (P) {\(v_7\)};
    \node[right] at (Q) {\(v_5\)};
    \node[font=\bfseries] at (0,2.65) {face \(234\)};
  \end{scope}
\end{tikzpicture}
\caption{The four three-dimensional cones of the fan of \(\PP^3\),
shown through their triangular cross-sections. Three faces are
subdivided by the new rays \(v_5,v_6,v_7\). The small triangles are
exactly the maximal cones listed above.}
\label{fig:fulton-nonprojective-threefold-fan}
\end{figure}

\clearpage
\begin{enumerate}[label=\textbf{Hint \arabic*.},wide=0pt]

\item \textbf{Begin with the fan of \(\PP^3\).}

Show that
\[
v_1+v_2+v_3+v_4=0
\]
and that the cones generated by any three of
\(v_1,v_2,v_3,v_4\) form the fan of \(\PP^3\).

Verify that
\[
v_5\in\Cone{v_3,v_4},
\qquad
v_6\in\Cone{v_1,v_4},
\qquad
v_7\in\Cone{v_2,v_4}.
\]
Use \cref{fig:fulton-nonprojective-threefold-fan} to check that the
listed cones subdivide the four maximal cones of the \(\PP^3\) fan.
Conclude that
\[
|\Sigma|=N_\RR,
\]
so \(X\) is complete.

\item \textbf{The example is smooth.}

For every maximal cone \(\sigma_{ijk}\), compute
\[
\det(v_i,v_j,v_k).
\]
Show that each determinant is \(\pm1\). Hence every maximal cone is
smooth and \(X\) is a smooth complete toric threefold.

\item \textbf{Translate projectivity into a support-function problem.}

Assume, for contradiction, that \(X\) is projective. Explain why there
must exist an integral strictly convex support function
\[
\psi:N_\RR\longrightarrow\RR.
\]
Set
\[
h_i=\psi(v_i).
\]

If \(\tau\) is a maximal cone with linear piece \(m_\tau\), recall
that strict convexity gives
\[
\langle m_\tau,u\rangle>\psi(u)
\qquad
\text{whenever }u\notin\tau.
\]

\item \textbf{The first incompatible inequality.}

Check the vector relation
\[
v_1+v_7=e_3=v_2+v_6.
\]
The vector \(e_3\) lies in the common face
\[
\Cone{v_2,v_6}
=
\sigma_{126}\cap\sigma_{267}.
\]
Use the linear piece on \(\sigma_{267}\) to show
\[
\psi(e_3)>h_1+h_7,
\]
and use the linear piece on \(\sigma_{126}\) to show
\[
\psi(e_3)=h_2+h_6.
\]
Deduce
\[
\boxed{h_2+h_6>h_1+h_7.}
\]

\item \textbf{Complete the cycle.}

Use the relations
\[
v_1+v_5=e_2=v_3+v_6
\]
and
\[
v_3+v_7=e_1=v_2+v_5
\]
together with the cones \(\sigma_{156}\) and \(\sigma_{357}\) to
obtain
\[
\boxed{h_1+h_5>h_3+h_6}
\]
and
\[
\boxed{h_3+h_7>h_2+h_5.}
\]

\item \textbf{Add the inequalities.}

Add the three boxed inequalities. Explain why the left- and right-hand
sides contain exactly the same six terms. Conclude that no strictly
convex support function can exist on \(\Sigma\).

Therefore \(X\) has no ample divisor and is not projective.

\item \textbf{Complete does not mean devoid of projective maps.}

The fan \(\Sigma\) refines the fan of \(\PP^3\). Explain why this gives
a proper birational toric morphism
\[
\pi:X\longrightarrow\PP^3.
\]
The pullback
\[
\pi^*\mathcal O_{\PP^3}(1)
\]
is basepoint free. Describe its support function on \(\Sigma\), and
explain why it is convex but not strictly convex.

Thus \(X\) has a nonconstant projective morphism even though \(X\)
itself is not projective.

\item \textbf{The missing polytope.}

Suppose there were a lattice polytope \(P\subseteq M_\RR\) with
\[
\Sigma_P=\Sigma.
\]
What support function would \(P\) produce? Use the contradiction above
to show that no such polytope exists.

Compare this with the surface theorem that every complete toric
surface fan is the normal fan of a lattice polygon. What genuinely new
phenomenon appears in dimension three?

\end{enumerate}

\begin{tcolorbox}[
  colback=PredictionTeal!5,
  colframe=PredictionTeal,
  title={The lesson of the example},
  fonttitle=\bfseries
]
Completeness asks only that the cones cover \(N_\RR\). Projectivity
asks for a global convex shape whose normal cones are exactly those
cones. In dimension three, a complete fan may contain a cyclic pattern
of walls whose required bends are mutually incompatible.
\end{tcolorbox}

\noindent
\textit{Reference.}
This is the classical example in Fulton, \emph{Introduction to Toric
Varieties}, \S3.4, p.~71. The strict-convexity contradiction is also
written out in Nicolas Perrin's lecture notes, Proposition~4.8.2.

\clearpage
\section*{Notes and references}
\addcontentsline{toc}{section}{Notes and references}
\markboth{Notes and references}{Notes and references}

The global-section formula and divisor polyhedra are developed in
Cox--Little--Schenck, \S4.3. Convexity, basepoint freeness, ampleness,
and normal fans are treated in Chapter~6. Projective toric varieties
from finite sets of lattice points are developed in Chapter~7.
Intersection numbers, volumes, and the surface polygon dictionary are
treated in Chapters~10 and~13.

The challenge exercise is Fulton's classical smooth complete
nonprojective toric threefold. It marks the precise point at which the
surface story stops generalizing: complete toric surfaces are
projective, but a complete three-dimensional fan need not admit any
strictly convex integral support function.

The next chapter introduces the Cox ring
\[
S=\Bbbk[x_\rho:\rho\in\Sigma(1)],
\]
which places one homogeneous coordinate on every invariant boundary
divisor and recovers the toric variety as a quotient of an open subset
of affine space.
